\documentclass[11pt,a4paper,openany]{book}
\usepackage[french]{babel}
\usepackage[utf8]{inputenc}
\usepackage[T1]{fontenc}
\usepackage{amsmath}
\usepackage{amsthm}
\usepackage{amsfonts}
\usepackage{amssymb}
\usepackage{graphicx}
\usepackage{thmtools}
\usepackage{verbatim}
\usepackage{pstricks-add}
\usepackage{pst-eucl}
\usepackage{tikz,tkz-tab}
\usepackage[left=1cm, right=1cm, top=0.5cm, bottom=0.5cm]{geometry}
\usepackage{hyperref}
\usepackage{url}
\usepackage{tcolorbox}
\tcbuselibrary{breakable}
\newtcolorbox{enoncebox}{breakable,colback=gray!6,colframe=black!30,title=Énoncé,fonttitle=\bfseries}
\newtcolorbox{methodebox}[1]{breakable,colback=blue!4,colframe=blue!40,title=#1,fonttitle=\bfseries}
\newtcolorbox{capacitesbox}{breakable,colback=orange!4,colframe=orange!70!brown,
  title={\textbf{Capacités attendues} \textemdash{} Programme officiel (BO)},
  fonttitle=\bfseries,before upper={\setlength\itemsep{2pt}}}
\newtcolorbox{demoprogbox}{breakable,colback=teal!4,colframe=teal!60!black,
  title={\textbf{Démonstrations exigibles} \textemdash{} Programme officiel (BO)},
  fonttitle=\bfseries,before upper={\setlength\itemsep{2pt}}}

\author{Ulrich TCHISSAMBOU}
\title{Mon livre de mathématiques de Terminale Spécialité}

\theoremstyle{definition}
\declaretheorem[thmbox=M, shaded={bgcolor={rgb}{2,1,1}}]{Définition}
\declaretheorem[thmbox=L]{Propriété}
\declaretheorem[thmbox=S]{Exemple}
\declaretheorem[thmbox=L]{Théorème}
\declaretheorem[thmbox=S]{Remarque}

\newcommand{\oijk}{(O;\vc{i},\vc{j},\vc{k})}
\newcommand{\vc}[1]{\overrightarrow{#1}}
\newcommand{\vca}[2]{
    \left(\begin{array}{c}
        #1\\
        #2
    \end{array}\right)}
\newcommand{\vcb}[3]{
    \left(\begin{array}{c}
        #1\\
        #2\\
        #3
    \end{array}\right)}
\newcommand{\nchap}{
    \setcounter{Définition}{0}
    \setcounter{Propriété}{0}
    \setcounter{Théorème}{0}
    \setcounter{Remarque}{0}
    \setcounter{Exemple}{0}
}

\begin{document}
    \renewcommand{\thesection}{\Alph{section})}
    \renewcommand{\thesubsection}{\arabic{subsection})}
    \tableofcontents

    % ============================================================
    \chapter*{Programme officiel --- Terminale Spécialité Mathématiques}
    \addcontentsline{toc}{chapter}{Programme officiel}
    % ============================================================

    Ce chapitre récapitule, conformément au \textbf{Bulletin officiel} (BO spécial
    n\degre{} 8 du 25 juillet 2019), l'ensemble des \textbf{capacités attendues}
    et des \textbf{démonstrations exigibles} pour chaque thème du programme de
    Terminale Spécialité Mathématiques.

    \bigskip

    % --- Thème 1 ------------------------------------------------
    \section*{Thème 1 \textemdash{} Récurrence et suites}
    \begin{capacitesbox}
    \begin{itemize}
        \item Mener un raisonnement par récurrence pour établir une propriété
              dépendant d'un entier $n$.
        \item Définir et utiliser les suites arithmétiques et géométriques
              (terme général, sommes).
        \item Étudier la monotonie d'une suite (signe de $u_{n+1}-u_n$,
              quotient, dérivée).
        \item Démontrer qu'une suite est majorée, minorée, bornée.
    \end{itemize}
    \end{capacitesbox}
    \begin{demoprogbox}
    \begin{itemize}
        \item $\displaystyle\sum_{k=1}^{n}k=\dfrac{n(n+1)}{2}$ par récurrence.
        \item Inégalité de Bernoulli : $(1+x)^n\geq 1+nx$
              pour $x>-1$, $n\in\mathbf{N}$.
    \end{itemize}
    \end{demoprogbox}

    % --- Thème 2 ------------------------------------------------
    \section*{Thème 2 \textemdash{} Limites de suites}
    \begin{capacitesbox}
    \begin{itemize}
        \item Calculer ou déterminer la limite d'une suite par les théorèmes
              d'opérations.
        \item Appliquer les théorèmes de comparaison (encadrement, suites
              adjacentes).
        \item Utiliser le théorème des suites monotones bornées (admis).
        \item Déterminer la limite d'une suite géométrique de raison $q$.
        \item Démontrer la convergence d'une suite définie par récurrence.
    \end{itemize}
    \end{capacitesbox}
    \begin{demoprogbox}
    \begin{itemize}
        \item Unicité de la limite d'une suite convergente (par l'absurde).
        \item Toute suite croissante et non majorée tend vers $+\infty$.
        \item $\displaystyle\lim_{n\to+\infty}q^n=0$ si $|q|<1$.
    \end{itemize}
    \end{demoprogbox}

    % --- Thème 3 ------------------------------------------------
    \section*{Thème 3 \textemdash{} Vecteurs, droites et plans de l'espace}
    \begin{capacitesbox}
    \begin{itemize}
        \item Représenter droites et plans dans l'espace ; décrire leurs
              positions relatives.
        \item Utiliser les vecteurs pour caractériser le parallélisme et la
              coplanarité.
        \item Déterminer une représentation paramétrique d'une droite ou d'un
              plan.
        \item Se repérer dans un repère orthonormé de l'espace.
    \end{itemize}
    \end{capacitesbox}
    \begin{demoprogbox}
    \begin{itemize}
        \item Si deux plans sécants contiennent chacun une droite de même
              direction, leur droite d'intersection est parallèle à ces droites.
        \item Une droite parallèle à deux droites sécantes d'un plan est
              parallèle à ce plan.
    \end{itemize}
    \end{demoprogbox}

    % --- Thème 4 ------------------------------------------------
    \section*{Thème 4 \textemdash{} Limites de fonctions}
    \begin{capacitesbox}
    \begin{itemize}
        \item Calculer des limites à l'aide des théorèmes d'opérations et de
              composition.
        \item Lever des formes indéterminées ($\infty-\infty$, $0/0$,
              $\infty/\infty$, $0\times\infty$).
        \item Utiliser les théorèmes de comparaison et les croissances
              comparées.
        \item Déterminer et interpréter les asymptotes.
    \end{itemize}
    \end{capacitesbox}
    \begin{demoprogbox}
    \begin{itemize}
        \item $\displaystyle\lim_{x\to 0}\dfrac{\sin x}{x}=1$
              (encadrement géométrique).
        \item Unicité de la limite finie d'une fonction en un point
              (par l'absurde).
    \end{itemize}
    \end{demoprogbox}

    % --- Thème 5 ------------------------------------------------
    \section*{Thème 5 \textemdash{} Dérivation et convexité}
    \begin{capacitesbox}
    \begin{itemize}
        \item Calculer la dérivée d'une fonction composée $v\circ u$ :
              $(v\circ u)'(x)=u'(x)\,v'(u(x))$.
        \item Caractériser la convexité par le signe de $f''$.
        \item Exploiter la convexité : comparaison $f(x)$ / tangente.
        \item Déterminer les points d'inflexion d'une courbe.
    \end{itemize}
    \end{capacitesbox}
    \begin{demoprogbox}
    \begin{itemize}
        \item Dérivation en chaîne : $[v(u(x))]'=u'(x)\,v'(u(x))$.
        \item Si $f''\geq 0$ sur $I$, la courbe est au-dessus de chacune de
              ses tangentes.
    \end{itemize}
    \end{demoprogbox}

    % --- Thème 6 ------------------------------------------------
    \section*{Thème 6 \textemdash{} Continuité}
    \begin{capacitesbox}
    \begin{itemize}
        \item Utiliser la continuité des fonctions usuelles et de leurs
              composées.
        \item Appliquer le théorème des valeurs intermédiaires (TVI) et son
              corollaire.
        \item Encadrer une solution de $f(x)=k$ par la méthode de dichotomie.
        \item Étudier la convergence d'une suite récurrente $u_{n+1}=f(u_n)$
              (point fixe).
    \end{itemize}
    \end{capacitesbox}
    \begin{demoprogbox}
    \begin{itemize}
        \item Cas particulier du TVI : $f$ continue sur $[a;b]$, $f(a)<0$,
              $f(b)>0 \Rightarrow \exists\, c\in{]a;b[}$ tel que $f(c)=0$.
        \item La dichotomie encadre la solution avec précision
              $\dfrac{b-a}{2^n}$ à l'étape $n$.
    \end{itemize}
    \end{demoprogbox}

    % --- Thème 7 ------------------------------------------------
    \section*{Thème 7 \textemdash{} Orthogonalité dans l'espace}
    \begin{capacitesbox}
    \begin{itemize}
        \item Calculer un produit scalaire dans l'espace (coordonnées ou
              normes).
        \item Démontrer l'orthogonalité de vecteurs, droites, droite et plan.
        \item Déterminer une équation cartésienne d'un plan à partir d'un
              vecteur normal.
        \item Calculer la distance d'un point à un plan.
    \end{itemize}
    \end{capacitesbox}
    \begin{demoprogbox}
    \begin{itemize}
        \item $\vc{u}\cdot\vc{v}=xx'+yy'+zz'$ dans un repère orthonormé.
        \item Distance d'un point $A(x_0,y_0,z_0)$ au plan $ax+by+cz+d=0$ :
              $d=\dfrac{|ax_0+by_0+cz_0+d|}{\sqrt{a^2+b^2+c^2}}$.
    \end{itemize}
    \end{demoprogbox}

    % --- Thème 8 ------------------------------------------------
    \section*{Thème 8 \textemdash{} Loi binomiale}
    \begin{capacitesbox}
    \begin{itemize}
        \item Identifier une expérience de Bernoulli et modéliser par
              $\mathcal{B}(n;p)$.
        \item Calculer $P(X=k)$, $P(X\leq k)$, $P(X\geq k)$.
        \item Calculer l'espérance, la variance et l'écart-type de
              $X\sim\mathcal{B}(n;p)$.
    \end{itemize}
    \end{capacitesbox}
    \begin{demoprogbox}
    \begin{itemize}
        \item Nombre de parties d'un ensemble à $n$ éléments : $2^n$
              (bijection avec les mots binaires).
        \item $P(X=k)=\dbinom{n}{k}p^k(1-p)^{n-k}$ (chemins dans un arbre).
    \end{itemize}
    \end{demoprogbox}

    % --- Thème 9 ------------------------------------------------
    \section*{Thème 9 \textemdash{} Logarithme népérien}
    \begin{capacitesbox}
    \begin{itemize}
        \item Utiliser les propriétés algébriques de $\ln$ :
              $\ln(ab)$, $\ln(a/b)$, $\ln(a^n)$.
        \item Résoudre équations et inéquations avec $\ln$.
        \item Étudier des fonctions contenant $\ln$ (variations, asymptotes).
        \item Croissances comparées : $\ln x \ll x^\alpha$ pour $\alpha>0$.
    \end{itemize}
    \end{capacitesbox}
    \begin{demoprogbox}
    \begin{itemize}
        \item $(\ln x)'=\dfrac{1}{x}$ (dérivée de la fonction réciproque de
              $\exp$).
        \item $\displaystyle\lim_{x\to+\infty}\dfrac{\ln x}{x}=0$.
        \item $\displaystyle\lim_{x\to 0^+}x\ln x=0$.
    \end{itemize}
    \end{demoprogbox}

    % --- Thème 10 -----------------------------------------------
    \section*{Thème 10 \textemdash{} Primitives et équations différentielles}
    \begin{capacitesbox}
    \begin{itemize}
        \item Calculer des primitives des fonctions usuelles.
        \item Résoudre $y'=ay$ et $y'=ay+b$ ($a\neq 0$).
        \item Résoudre un problème de Cauchy (condition initiale).
        \item Modéliser un phénomène par une équation différentielle.
    \end{itemize}
    \end{capacitesbox}
    \begin{demoprogbox}
    \begin{itemize}
        \item Les solutions de $y'=ay$ sont exactement les
              $x\mapsto Ce^{ax}$, $C\in\mathbf{R}$.
        \item Les solutions de $y'=ay+b$ sont les
              $x\mapsto Ce^{ax}-\dfrac{b}{a}$, $C\in\mathbf{R}$.
    \end{itemize}
    \end{demoprogbox}

    % --- Thème 11 -----------------------------------------------
    \section*{Thème 11 \textemdash{} Combinatoire et dénombrement}
    \begin{capacitesbox}
    \begin{itemize}
        \item Appliquer les principes additif et multiplicatif.
        \item Calculer permutations $n!$, arrangements $A_n^k$, combinaisons
              $\dbinom{n}{k}$.
        \item Utiliser la formule de Pascal et développer $(a+b)^n$.
    \end{itemize}
    \end{capacitesbox}
    \begin{demoprogbox}
    \begin{itemize}
        \item Nombre de parties d'un ensemble à $n$ éléments : $2^n$.
        \item Formule de Pascal :
              $\dbinom{n}{k}+\dbinom{n}{k+1}=\dbinom{n+1}{k+1}$.
    \end{itemize}
    \end{demoprogbox}

    % --- Thème 12 -----------------------------------------------
    \section*{Thème 12 \textemdash{} Fonctions cosinus et sinus}
    \begin{capacitesbox}
    \begin{itemize}
        \item Connaître les valeurs exactes pour les angles remarquables.
        \item Utiliser les formules d'addition et de double-angle.
        \item Résoudre $\cos x=k$, $\sin x=k$, $\cos x=\cos a$,
              $\sin x=\sin a$.
        \item Étudier des fonctions faisant intervenir $\cos$ ou $\sin$.
    \end{itemize}
    \end{capacitesbox}
    \begin{demoprogbox}
    \begin{itemize}
        \item Formules d'addition :
              $\cos(a+b)=\cos a\cos b-\sin a\sin b$ et
              $\sin(a+b)=\sin a\cos b+\cos a\sin b$.
        \item $(\cos x)'=-\sin x$ et $(\sin x)'=\cos x$.
    \end{itemize}
    \end{demoprogbox}

    % --- Thème 13 -----------------------------------------------
    \section*{Thème 13 \textemdash{} Calcul intégral}
    \begin{capacitesbox}
    \begin{itemize}
        \item Calculer $\displaystyle\int_a^b f(x)\,\mathrm{d}x=[F(x)]_a^b$.
        \item Interpréter l'intégrale comme une aire algébrique.
        \item Calculer une valeur moyenne.
        \item Intégration par parties :
              $\displaystyle\int_a^b uv'\,\mathrm{d}x=[uv]_a^b
              -\int_a^b u'v\,\mathrm{d}x$.
    \end{itemize}
    \end{capacitesbox}
    \begin{demoprogbox}
    \begin{itemize}
        \item Si $f\geq 0$ sur $[a;b]$, alors
              $\displaystyle\int_a^b f(x)\,\mathrm{d}x\geq 0$.
        \item Inégalité de la moyenne :
              $\left|\displaystyle\int_a^b f\right|\leq(b-a)
              \displaystyle\max_{[a;b]}|f|$.
        \item Démonstration de la formule d'intégration par parties.
    \end{itemize}
    \end{demoprogbox}

    % --- Thème 14 -----------------------------------------------
    \section*{Thème 14 \textemdash{} Loi des grands nombres}
    \begin{capacitesbox}
    \begin{itemize}
        \item Calculer espérance, variance et écart-type d'une variable
              aléatoire.
        \item Comprendre et utiliser l'inégalité de Bienaymé-Tchebychev.
        \item Interpréter et appliquer la loi des grands nombres.
    \end{itemize}
    \end{capacitesbox}
    \begin{demoprogbox}
    \begin{itemize}
        \item $E(X)=np$ et $V(X)=np(1-p)$ pour $X\sim\mathcal{B}(n;p)$.
        \item Inégalité de Bienaymé-Tchebychev :
              $P(|X-E(X)|\geq\delta)\leq\dfrac{V(X)}{\delta^2}$.
        \item Loi faible des grands nombres (admise au lycée).
    \end{itemize}
    \end{demoprogbox}

    \newpage

    \chapter{Récurrence, Suites}

    \begin{capacitesbox}
    \begin{itemize}
        \item Mener un raisonnement par récurrence pour établir une propriété dépendant d'un entier $n$.
        \item Identifier et utiliser les suites arithmétiques et géométriques (terme général, somme).
        \item Étudier la monotonie d'une suite (signe de $u_{n+1}-u_n$, quotient, dérivée).
        \item Démontrer qu'une suite est majorée, minorée, bornée.
    \end{itemize}
    \end{capacitesbox}

    \begin{demoprogbox}
    \begin{itemize}
        \item Démontrer par récurrence que $\displaystyle\sum_{k=1}^{n}k=\dfrac{n(n+1)}{2}$.
        \item Démontrer l'inégalité de Bernoulli : pour tout $x>-1$ et $n\in\mathbf{N}$,
              $(1+x)^n\geq 1+nx$.
    \end{itemize}
    \end{demoprogbox}

    \section{Récurrence}
    \begin{Définition}
        Une propriété mathématique est une phrase qui peut être vraie ou fausse.
    \end{Définition}

    \begin{Théorème}
        Soient $n_0$ un entier et $P(n)$ une propriété dépendant de l'entier naturel $n$
        définie pour $n\geq n_0$.\\
        On suppose :
        \begin{enumerate}
            \item \textbf{Initialisation :} $P(n_0)$ est vraie.
            \item \textbf{Hérédité :} Pour tout entier naturel $k\geq n_0$ fixé, si $P(k)$
                  est vraie alors $P(k+1)$ est vraie.
        \end{enumerate}
        Alors pour tout entier naturel $n\geq n_0$, $P(n)$ est vraie.
    \end{Théorème}

    \begin{Exemple}
        Soit $(u_n)$ une suite définie par $u_0=3$ et pour tout entier naturel $n$ par
        $u_{n+1}=3u_n-2$.
        Démontrer par récurrence que, pour tout entier naturel $n$, $u_n=2\times 3^n+1$.\\
        Notons $P(n) : \ll u_n=2\times 3^n+1\gg$.
        \begin{enumerate}
            \item \textbf{Initialisation :} Pour $n=0$ on a $2\times 3^0+1=3$,
                  donc $P(0)$ est vraie.
            \item \textbf{Hérédité :} Supposons la propriété vraie pour un $k\geq 0$
                  (i.e.\ $u_k=2\times 3^k+1$). On veut montrer que $P(k+1)$
                  ($u_{k+1}=2\times 3^{k+1}+1$) est vraie.
            \begin{align*}
                u_{k+1}&=3u_k-2\\
                       &=3\times (2\times 3^k+1)-2\\
                       &=2\times 3\times 3^k+3-2\\
                       &=2\times 3^{k+1}+1
            \end{align*}
            Donc $P(k+1)$ est vraie.
            \item \textbf{Conclusion :} Comme $P(0)$ est vraie et que pour tout entier $k$,
                  $P(k)$ vraie implique $P(k+1)$ vraie, la propriété $P(n)$ est vraie pour
                  tout $n\in\mathbf{N}$.
        \end{enumerate}
    \end{Exemple}

    \begin{Exemple}
    Démontrons par récurrence que, pour tout entier naturel $n\geq 1$, la somme des $n$
    premiers entiers est :
    \[
        1 + 2 + \cdots + n = \frac{n(n+1)}{2}.
    \]
    \begin{enumerate}
        \item \textbf{Initialisation :} Pour $n = 1$ :
        \[
            1 = \frac{1(1+1)}{2} = \frac{2}{2} = 1.
        \]
        La propriété est vraie pour $n = 1$.
        \item \textbf{Hérédité :} Supposons la propriété vraie pour un entier $k \geq 1$,
        c'est-à-dire :
        \[
            1 + 2 + \cdots + k = \frac{k(k+1)}{2}.
        \]
        Montrons qu'elle est vraie pour $k+1$ :
        \begin{align*}
            1 + 2 + \cdots + k + (k+1)
            &= \frac{k(k+1)}{2} + (k+1) \\
            &= \frac{k(k+1) + 2(k+1)}{2} \\
            &= \frac{(k+1)(k+2)}{2}.
        \end{align*}
        La propriété est vraie pour $k+1$.
        \item \textbf{Conclusion :} La propriété est vraie pour $n=1$ et héréditaire,
        donc vraie pour tout entier naturel $n\geq 1$.
    \end{enumerate}
    \end{Exemple}

    \begin{Exemple}[Contre-exemple — importance de l'initialisation]
    Montrons que la proposition « pour tout entier naturel $n \geq 1$, $2^n \geq n^2$ »
    est \textbf{fausse} en exhibant un contre-exemple.
    \begin{itemize}
        \item Pour $n = 1$ : $2^1 = 2 \geq 1 = 1^2$. \quad Vraie.
        \item Pour $n = 2$ : $2^2 = 4 = 2^2$. \quad Vraie.
        \item Pour $n = 3$ : $2^3 = 8 < 9 = 3^2$. \quad \textbf{Fausse.}
    \end{itemize}
    La propriété est donc fausse dès $n=3$.\\
    \textbf{Remarque pédagogique :} même si l'hérédité est vérifiable pour $k\geq 3$,
    l'échec de l'initialisation en $n=3$ suffit à invalider la proposition.
    Cet exemple illustre l'importance de vérifier rigoureusement l'initialisation.
    \end{Exemple}

    \section{Monotonie d'une suite}
    \begin{Définition}
        Soit $(u_n)$ une suite de nombres réels.
        \begin{itemize}
            \item La suite $(u_n)$ est \textbf{croissante} si pour tout $n\in\mathbf{N}$,
                  $u_n\leq u_{n+1}$.
            \item La suite $(u_n)$ est \textbf{décroissante} si pour tout $n\in\mathbf{N}$,
                  $u_n\geq u_{n+1}$.
            \item La suite $(u_n)$ est \textbf{monotone} si elle est croissante ou décroissante.
        \end{itemize}
    \end{Définition}

    \subsection{Comment montrer qu'une suite est monotone ?}

    \subsubsection{Méthode 1 — Signe de $u_{n+1}-u_n$}
    On étudie le signe de $u_{n+1}-u_n$.
    \begin{itemize}
        \item $u_{n+1}-u_n \geq 0$ pour tout $n$ $\Rightarrow$ suite croissante.
        \item $u_{n+1}-u_n \leq 0$ pour tout $n$ $\Rightarrow$ suite décroissante.
        \item $u_{n+1}-u_n > 0$ pour tout $n$ $\Rightarrow$ suite strictement croissante.
        \item $u_{n+1}-u_n < 0$ pour tout $n$ $\Rightarrow$ suite strictement décroissante.
    \end{itemize}

    \subsubsection{Méthode 2 — Quotient $\dfrac{u_{n+1}}{u_n}$ (termes strictement positifs)}
    Si $(u_n)$ est strictement positive :
    \begin{itemize}
        \item $\dfrac{u_{n+1}}{u_n} \geq 1$ pour tout $n$ $\Rightarrow$ suite croissante.
        \item $\dfrac{u_{n+1}}{u_n} \leq 1$ pour tout $n$ $\Rightarrow$ suite décroissante.
        \item $\dfrac{u_{n+1}}{u_n} > 1$ pour tout $n$ $\Rightarrow$ suite strictement croissante.
        \item $\dfrac{u_{n+1}}{u_n} < 1$ pour tout $n$ $\Rightarrow$ suite strictement décroissante.
    \end{itemize}

    \subsubsection{Méthode 3 — Sens de variation d'une fonction}
    S'il existe une fonction $f$ telle que $u_n=f(n)$, on étudie le sens de variation de $f$
    sur $[0;+\infty[$ via la dérivée $f'$.
    \begin{itemize}
        \item $f'(x) \geq 0$ sur $[0;+\infty[$ $\Rightarrow$ $(u_n)$ croissante.
        \item $f'(x) \leq 0$ sur $[0;+\infty[$ $\Rightarrow$ $(u_n)$ décroissante.
        \item $f'(x) > 0$ sur $[0;+\infty[$ $\Rightarrow$ $(u_n)$ strictement croissante.
        \item $f'(x) < 0$ sur $[0;+\infty[$ $\Rightarrow$ $(u_n)$ strictement décroissante.
    \end{itemize}

    \subsubsection{Méthode 4 — Récurrence sur $P(n) : u_n \leq u_{n+1}$}
    \begin{itemize}
        \item \textbf{Initialisation :} vérifier $u_0 \leq u_1$ (ou $u_1 \leq u_2$).
        \item \textbf{Hérédité :} supposer $u_k \leq u_{k+1}$ et montrer $u_{k+1} \leq u_{k+2}$.
        \item Si les deux étapes sont vérifiées, la suite est croissante.
              Si on montre $u_n \geq u_{n+1}$ par récurrence, elle est décroissante.
    \end{itemize}

    \begin{Définition}
        Soit $(u_n)$ une suite de nombres réels.
        \begin{itemize}
            \item On dit que $(u_n)$ est \textbf{majorée} s'il existe un réel $M$ tel que
                  pour tout $n\in\mathbf{N}$, $u_n\leq M$.
            \item On dit que $(u_n)$ est \textbf{minorée} s'il existe un réel $m$ tel que
                  pour tout $n\in\mathbf{N}$, $u_n\geq m$.
            \item On dit que $(u_n)$ est \textbf{bornée} si elle est minorée et majorée.
        \end{itemize}
    \end{Définition}

    \chapter{Limites de suite}
    \nchap

    \begin{capacitesbox}
    \begin{itemize}
        \item Calculer ou conjecturer la limite d'une suite à l'aide des théorèmes d'opérations.
        \item Appliquer les théorèmes de comparaison (encadrement, suites adjacentes).
        \item Utiliser le théorème des suites monotones bornées (admis).
        \item Déterminer la limite d'une suite géométrique de raison $q$.
        \item Démontrer la convergence d'une suite définie par récurrence.
    \end{itemize}
    \end{capacitesbox}

    \begin{demoprogbox}
    \begin{itemize}
        \item Démontrer l'unicité de la limite d'une suite convergente (raisonnement par l'absurde).
        \item Démontrer que toute suite croissante et non majorée tend vers $+\infty$.
        \item Démontrer que $\displaystyle\lim_{n\to+\infty}q^n=0$ si $|q|<1$.
    \end{itemize}
    \end{demoprogbox}

    \section{Limite}
    \subsection{Limite finie}
    \begin{Définition}
    Une suite $(u_n)$ converge vers le réel $\ell$ si tout intervalle contenant $\ell$
    contient tous les termes de la suite à partir d'un certain rang.
    On écrit $\lim\limits_{n\to+\infty}u_n=\ell$.
    On dit que la suite $(u_n)$ est \textbf{convergente} et converge vers $\ell$.
    \end{Définition}

    \begin{Remarque}
    Une suite $(u_n)$ converge vers $\ell$ si, pour tout $\varepsilon>0$, l'intervalle
    $]\ell-\varepsilon;\ell+\varepsilon[$ contient tous les termes de la suite à partir
    d'un certain rang.
    \end{Remarque}

    \begin{center}
    \begin{tikzpicture}[scale=0.9]
        \draw[->] (0,0) -- (7,0) node[right] {$n$};
        \draw[->] (0,0) -- (0,3) node[above] {$u_n$};
        \foreach \n/\y in {1/2.5,2/1.8,3/1.4,4/1.2,5/1.1,6/1.05}
            \fill (\n,\y) circle (2pt);
        \draw[dashed,thick] (0,1) -- (7,1) node[right] {$\ell$};
        \draw[dashed,red] (0,0.9) -- (7,0.9);
        \draw[dashed,red] (0,1.1) -- (7,1.1);
        \node[red] at (7.7,1) {$\ell \pm \varepsilon$};
    \end{tikzpicture}
    \end{center}

    \begin{Définition}
    Une suite qui n'est pas convergente est dite \textbf{divergente}.
    \end{Définition}

    \begin{Exemple}
    $u_n=(-1)^n$ est une suite divergente ; elle n'admet pas de limite lorsque $n$ tend
    vers l'infini.
    \end{Exemple}

    \begin{center}
    \begin{tikzpicture}[scale=0.9]
        \draw[->] (0,0) -- (7,0) node[right] {$n$};
        \draw[->] (0,-1.5) -- (0,1.5) node[above] {$u_n$};
        \foreach \n/\y in {1/1,2/-1,3/1,4/-1,5/1,6/-1}
            \fill (\n,\y) circle (2pt);
        \draw[dashed,thick] (0,1) -- (7,1);
        \draw[dashed,thick] (0,-1) -- (7,-1);
    \end{tikzpicture}
    \end{center}

    \begin{Propriété}[Unicité de la limite]
    Si une suite $(u_n)$ a une limite réelle, alors cette limite est unique.
    \end{Propriété}
    \begin{proof}
    Supposons que $(u_n)$ ait deux limites $\ell_1$ et $\ell_2$ avec $\ell_1 \neq \ell_2$.
    Posons $\varepsilon = \dfrac{|\ell_1 - \ell_2|}{3} > 0$.
    Par définition de la limite, il existe $N_1$ tel que pour tout $n \geq N_1$,
    $|u_n - \ell_1| < \varepsilon$, et il existe $N_2$ tel que pour tout $n \geq N_2$,
    $|u_n - \ell_2| < \varepsilon$.
    Donc, pour $n \geq \max(N_1,N_2)$ :
    \[
        |\ell_1 - \ell_2| \leq |u_n - \ell_1| + |u_n - \ell_2| < 2\varepsilon
        = \tfrac{2}{3}|\ell_1 - \ell_2|.
    \]
    Ceci est impossible si $\ell_1 \neq \ell_2$, donc $\ell_1 = \ell_2$.
    \end{proof}

    \subsection{Quelques limites finies classiques}
    \begin{Propriété}
    \begin{enumerate}
        \item $\lim\limits_{n\to+\infty}\dfrac{1}{\sqrt{n}}=0$
        \item Pour tout $k\in\mathbf{N}^*$, $\lim\limits_{n\to+\infty}\dfrac{1}{n^k}=0$
        \item Pour tout $-1<q<1$, $\lim\limits_{n\to+\infty}q^n=0$
    \end{enumerate}
    \end{Propriété}

    \begin{Exemple}[Application du point 1 : limite de $\dfrac{1}{\sqrt{n}}$]
    Soit $u_n=\dfrac{3}{\sqrt{n}}$.
    \[
        \lim_{n\to+\infty}u_n
        = \lim_{n\to+\infty}\frac{3}{\sqrt{n}}
        = 3\times\lim_{n\to+\infty}\frac{1}{\sqrt{n}}
        = 3\times 0 = 0.
    \]
    De même, $v_n=\dfrac{\sqrt{n}+5}{\sqrt{n}}=1+\dfrac{5}{\sqrt{n}}\xrightarrow[n\to+\infty]{}1+0=1$.
    \end{Exemple}

    \begin{Exemple}[Application du point 2 : limite de $\dfrac{1}{n^k}$]
    \begin{enumerate}
        \item Pour $k=2$ : $\displaystyle\lim_{n\to+\infty}\frac{1}{n^2}=0$.
              Donc $\displaystyle\lim_{n\to+\infty}\frac{4n+1}{n^2}
              =\lim_{n\to+\infty}\frac{4}{n}+\frac{1}{n^2}=0$.
        \item Pour $k=3$ : $\displaystyle\lim_{n\to+\infty}\frac{2n^2-1}{n^3}
              =\lim_{n\to+\infty}\frac{2}{n}-\frac{1}{n^3}=0$.
        \item \textbf{Règle pratique :} dans une fraction rationnelle en $n$,
              on factorise par le terme dominant du dénominateur :
              \[
                  \frac{3n^2+n-2}{5n^2-4}
                  = \frac{n^2\!\left(3+\dfrac{1}{n}-\dfrac{2}{n^2}\right)}
                         {n^2\!\left(5-\dfrac{4}{n^2}\right)}
                  \xrightarrow[n\to+\infty]{}\frac{3}{5}.
              \]
    \end{enumerate}
    \end{Exemple}

    \begin{Exemple}[Application du point 3 : suite géométrique $q^n$ avec $|q|<1$]
    \begin{enumerate}
        \item $q=\dfrac{1}{2}$ : $\displaystyle\lim_{n\to+\infty}\left(\frac{1}{2}\right)^n=0$.
              Les termes valent $1,\,\frac{1}{2},\,\frac{1}{4},\,\frac{1}{8},\ldots$
        \item $q=-\dfrac{1}{3}$ : $|q|=\dfrac{1}{3}<1$, donc
              $\displaystyle\lim_{n\to+\infty}\left(-\frac{1}{3}\right)^n=0$
              (les termes oscillent mais se rapprochent de $0$).
        \item $q=0{,}99$ : même si $q$ est proche de $1$, comme $|q|<1$ on a
              $\displaystyle\lim_{n\to+\infty}(0{,}99)^n=0$.
        \item \textbf{Cas limite important :} si $q=1$, $q^n=1$ pour tout $n$ ;
              si $|q|>1$, $|q^n|\to+\infty$.
              Seul $|q|<1$ garantit $q^n\to 0$.
    \end{enumerate}
    \end{Exemple}

    \subsection{Limites infinies classiques}

    \begin{Propriété}
    \begin{enumerate}
        \item Pour tout $k\in\mathbf{N}^*$,
              $\lim\limits_{n\to+\infty}n^k=+\infty$
        \item Pour tout $q>1$,
              $\lim\limits_{n\to+\infty}q^n=+\infty$
        \item $\lim\limits_{n\to+\infty}n!=+\infty$
        \item $\lim\limits_{n\to+\infty}\sqrt{n}=+\infty$
    \end{enumerate}
    \end{Propriété}

    \begin{Exemple}[Limites de suites polynomiales et exponentielles]
    \begin{enumerate}
        \item $u_n=3n^2-5n+2$. Le terme dominant est $3n^2$, donc :
              \[
                  \lim_{n\to+\infty}u_n
                  =\lim_{n\to+\infty}n^2\!\left(3-\frac{5}{n}+\frac{2}{n^2}\right)
                  =+\infty.
              \]
        \item $v_n=2^n-100n^3$. L'exponentielle l'emporte sur le polynôme :
              \[
                  v_n=2^n\!\left(1-\frac{100n^3}{2^n}\right)
                  \xrightarrow[n\to+\infty]{}+\infty.
              \]
        \item $w_n=(-2)^n$ : comme $|{-2}|=2>1$, $|w_n|=2^n\to+\infty$,
              mais la suite \textbf{diverge} (elle n'a pas de limite car elle oscille).
    \end{enumerate}
    \end{Exemple}

    \subsection{Croissances comparées}

    \begin{Propriété}[Hiérarchie des suites classiques]
    Pour tout $q>1$ et tout $k\in\mathbf{N}^*$, on a l'ordre de croissance :
    \[
        \ln n \;\ll\; n^k \;\ll\; q^n \;\ll\; n!
    \]
    ce qui signifie concrètement :
    \begin{enumerate}
        \item $\displaystyle\lim_{n\to+\infty}\frac{\ln n}{n^k}=0$
              \quad (le logarithme est négligeable devant toute puissance)
        \item $\displaystyle\lim_{n\to+\infty}\frac{n^k}{q^n}=0$
              \quad (toute puissance est négligeable devant toute exponentielle)
        \item $\displaystyle\lim_{n\to+\infty}\frac{q^n}{n!}=0$
              \quad (toute exponentielle est négligeable devant $n!$)
    \end{enumerate}
    \end{Propriété}

    \begin{Exemple}[Croissances comparées --- applications]
    \begin{enumerate}
        \item \textbf{Puissance vs exponentielle.}
              $\displaystyle\lim_{n\to+\infty}\frac{n^5}{2^n}=0$ :
              même si $n^5$ croit très vite, $2^n$ finit par l'emporter.

              \medskip
              On peut l'observer numériquement :
              \begin{center}
              \begin{tabular}{|c|r|r|r|}
                  \hline
                  $n$ & $n^5$ & $2^n$ & $\dfrac{n^5}{2^n}$ \\
                  \hline
                  10  & $100\,000$        & $1\,024$           & $97{,}7$ \\
                  20  & $3\,200\,000$     & $1\,048\,576$      & $3{,}05$ \\
                  50  & $312\,500\,000$   & $\approx 10^{15}$  & $\approx 0$ \\
                  \hline
              \end{tabular}
              \end{center}

        \item \textbf{Limite d'une suite mixte.}
              \[
                  u_n=\frac{3n^3+n}{5^n}\,.
              \]
              Comme $\dfrac{n^3}{5^n}\to 0$, on a
              $u_n=\dfrac{3n^3}{5^n}+\dfrac{n}{5^n}\to 0+0=0$.

        \item \textbf{Limite indéterminée levée par croissances comparées.}
              \[
                  v_n=\frac{2^n-n^{10}}{2^n+1}
                  =\frac{2^n\!\left(1-\dfrac{n^{10}}{2^n}\right)}
                        {2^n\!\left(1+\dfrac{1}{2^n}\right)}
                  \xrightarrow[n\to+\infty]{}\frac{1-0}{1+0}=1.
              \]

        \item \textbf{Comparaison $\ln n$ vs $\sqrt{n}$.}
              $\dfrac{\ln n}{\sqrt{n}}\to 0$ : le logarithme est négligeable
              même devant $\sqrt{n}$.
    \end{enumerate}
    \end{Exemple}

    \subsection{Suites récurrentes --- limites classiques}

    \begin{Propriété}[Suite arithmético-géométrique]
    Soit $(u_n)$ définie par $u_0\in\mathbf{R}$ et $u_{n+1}=au_n+b$
    avec $a\neq 1$. La limite de $(u_n)$ est :
    \begin{itemize}
        \item Si $|a|<1$ : $(u_n)$ converge vers le point fixe
              $\ell=\dfrac{b}{1-a}$.
        \item Si $|a|>1$ : $(u_n)$ diverge vers $\pm\infty$
              (sauf si $u_0=\ell$).
        \item Si $a=1$ : $(u_n)$ est arithmétique de raison $b$.
    \end{itemize}
    \end{Propriété}

    \begin{Exemple}[Suite arithmético-géométrique]
    Soit $(u_n)$ définie par $u_0=1$ et $u_{n+1}=\dfrac{1}{2}u_n+3$.
    \begin{enumerate}
        \item \textbf{Point fixe :} $\ell=\dfrac{1}{2}\ell+3
              \Rightarrow \ell=6$.
        \item \textbf{Suite auxiliaire :} $v_n=u_n-6$ vérifie
              $v_{n+1}=\dfrac{1}{2}v_n$, donc $v_n=v_0\times\!\left(\dfrac{1}{2}\right)^n
              =(1-6)\times\!\left(\dfrac{1}{2}\right)^n
              =-5\times\!\left(\dfrac{1}{2}\right)^n$.
        \item \textbf{Terme général :}
              $u_n=6-5\times\!\left(\dfrac{1}{2}\right)^n$.
        \item \textbf{Limite :} comme $\left|\dfrac{1}{2}\right|<1$,
              $\left(\dfrac{1}{2}\right)^n\to 0$, donc
              $\displaystyle\lim_{n\to+\infty}u_n=6=\ell$.
    \end{enumerate}
    \end{Exemple}

    \begin{Théorème}
    \begin{itemize}
        \item Une suite croissante et majorée est convergente.
        \item Une suite décroissante et minorée est convergente.
    \end{itemize}
    \end{Théorème}

    \subsection{Limite infinie}
    \begin{Définition}
    Une suite $(u_n)$ a pour limite $+\infty$ si, pour tout réel $A>0$, l'intervalle
    $[A;+\infty[$ contient tous les termes de la suite à partir d'un certain rang.
    On écrit $\lim\limits_{n\to+\infty}u_n=+\infty$.
    \end{Définition}

    \begin{center}
    \begin{tikzpicture}[scale=0.9]
        \draw[->] (0,0) -- (7,0) node[right] {$n$};
        \draw[->] (0,0) -- (0,4) node[above] {$u_n$};
        \foreach \n/\y in {1/0.5,2/1,3/1.8,4/2.5,5/3.2,6/3.8}
            \fill (\n,\y) circle (2pt);
    \end{tikzpicture}
    \end{center}

    \begin{Définition}
    Une suite $(u_n)$ a pour limite $-\infty$ si, pour tout réel $B<0$, l'intervalle
    $]-\infty;B]$ contient tous les termes de la suite à partir d'un certain rang.
    On écrit $\lim\limits_{n\to+\infty}u_n=-\infty$.
    \end{Définition}

    \section{Opérations sur les limites}
    \begin{Définition}
    Une \textbf{forme indéterminée} est une opération rencontrée lors du calcul d'une
    limite pour laquelle on ne peut pas conclure directement.
    \end{Définition}

    \subsection{Somme}
    \begin{Propriété}
    \begin{center}
    \begin{tabular}{|c|c|c|}
        \hline
        $\lim u_n$ & $\lim v_n$ & $\lim(u_n+v_n)$ \\
        \hline
        $\ell_1$ & $\ell_2$ & $\ell_1+\ell_2$ \\
        \hline
        $\ell_1$ & $+\infty$ & $+\infty$ \\
        \hline
        $\ell_1$ & $-\infty$ & $-\infty$ \\
        \hline
        $+\infty$ & $+\infty$ & $+\infty$ \\
        \hline
        $-\infty$ & $-\infty$ & $-\infty$ \\
        \hline
        $+\infty$ & $-\infty$ & \textbf{F.I.} \\
        \hline
    \end{tabular}
    \end{center}
    \end{Propriété}

    \subsection{Produit}
    \begin{Propriété}
    \begin{center}
    \begin{tabular}{|c|c|c|}
        \hline
        $\lim u_n$ & $\lim v_n$ & $\lim(u_n\times v_n)$ \\
        \hline
        $\ell_1$ & $\ell_2$ & $\ell_1\times\ell_2$ \\
        \hline
        $\ell_1>0$ & $+\infty$ & $+\infty$ \\
        \hline
        $\ell_1>0$ & $-\infty$ & $-\infty$ \\
        \hline
        $\ell_1<0$ & $+\infty$ & $-\infty$ \\
        \hline
        $\ell_1<0$ & $-\infty$ & $+\infty$ \\
        \hline
        $+\infty$ & $+\infty$ & $+\infty$ \\
        \hline
        $-\infty$ & $-\infty$ & $+\infty$ \\
        \hline
        $0$ & $\pm\infty$ & \textbf{F.I.} \\
        \hline
    \end{tabular}
    \end{center}
    \end{Propriété}

    \subsection{Quotient}
    \begin{Propriété}
    \begin{center}
    \begin{tabular}{|c|c|c|}
        \hline
        $\lim u_n$ & $\lim v_n$ & $\lim\dfrac{u_n}{v_n}$ \\
        \hline
        $\ell_1$ & $\ell_2\neq 0$ & $\dfrac{\ell_1}{\ell_2}$ \\
        \hline
        $\ell_1$ & $\pm\infty$ & $0$ \\
        \hline
        $\pm\infty$ & $\pm\infty$ & \textbf{F.I.} \\
        \hline
        $0$ & $0$ & \textbf{F.I.} \\
        \hline
        $\ell_1>0$ ou $+\infty$ & $0$ avec $v_n>0$ & $+\infty$ \\
        \hline
        $\ell_1<0$ ou $-\infty$ & $0$ avec $v_n>0$ & $-\infty$ \\
        \hline
        $\ell_1>0$ ou $+\infty$ & $0$ avec $v_n<0$ & $-\infty$ \\
        \hline
        $\ell_1<0$ ou $-\infty$ & $0$ avec $v_n<0$ & $+\infty$ \\
        \hline
    \end{tabular}
    \end{center}
    \end{Propriété}

    \begin{Exemple}
    \begin{enumerate}
        \item $\lim\limits_{n\to+\infty} n^2+3n-5=+\infty$
        \item $\lim\limits_{n\to+\infty} \left(3+\dfrac{2}{n^3}\right)\left(\dfrac{1}{n}+5\right)=15$
    \end{enumerate}
    \end{Exemple}

    \subsection{Comment lever une indétermination ?}
    On peut effectuer une ou plusieurs factorisations.
    \begin{Exemple}
    $\lim\limits_{n\to+\infty} n^3-2n^2+5=+\infty$ (après factorisation par $n^3$).
    \end{Exemple}

    \section{Limites et comparaison}
    \begin{Propriété}
    Si $(u_n)$ et $(v_n)$ sont deux suites telles que, à partir d'un certain rang,
    $v_n\leq u_n$ et $\lim v_n=+\infty$, alors $\lim u_n=+\infty$.
    \end{Propriété}

    \begin{Propriété}
    Si $(u_n)$ et $(v_n)$ sont deux suites telles que, à partir d'un certain rang,
    $u_n\leq v_n$ et $\lim v_n=-\infty$, alors $\lim u_n=-\infty$.
    \end{Propriété}

    \begin{Théorème}[Théorème des gendarmes]
    Si $(u_n)$, $(v_n)$ et $(w_n)$ sont trois suites telles que, à partir d'un certain rang,
    $v_n\leq u_n\leq w_n$ et $\lim v_n=\lim w_n=\ell\in\mathbf{R}$, alors $\lim u_n=\ell$.
    \end{Théorème}

    \section{Suites de référence}
    \begin{Propriété}
    Pour toute suite arithmétique $(u_n)$ de raison $r$ :
    \begin{itemize}
        \item Si $r>0$ alors $(u_n)$ diverge vers $+\infty$.
        \item Si $r<0$ alors $(u_n)$ diverge vers $-\infty$.
    \end{itemize}
    \end{Propriété}

    \begin{Propriété}
    Pour toute suite géométrique $(u_n)$ de raison $q$ :
    \begin{itemize}
        \item Si $q>1$ alors $(u_n)$ diverge vers $+\infty$.
        \item Si $q<-1$ alors $(u_n)$ diverge et n'a pas de limite.
        \item Si $-1<q<1$ alors $(u_n)$ converge vers $0$.
    \end{itemize}
    \end{Propriété}

    \section{Utilisation de la calculatrice}
    \subsection*{NumWorks}
    Menu \textbf{Suites} : entrer la suite en explicite ou récurrente, puis observer dans
    \textbf{Tableau} ou \textbf{Graphique}.
    \subsection*{TI}
    Mettre le mode \textbf{Suites}, définir $u(n)$, utiliser \textbf{TABLE} et \textbf{GRAPH}.
    \subsection*{Casio}
    Menu \textbf{RECUR}, entrer $u_n$, visualiser le tableau et le graphe.
    \medskip\\
    \textbf{Remarque :} Ces outils servent à \emph{conjecturer} une limite, mais pas à la
    prouver rigoureusement.

    \section{Algorithme de seuil d'une suite}
    \begin{Définition}
    On appelle \textbf{seuil} d'une suite $(u_n)$ par rapport à un réel $\ell$ et une
    précision $\varepsilon>0$ le plus petit entier $N$ tel que, pour tout $n \geq N$,
    $|u_n-\ell|<\varepsilon$.
    \end{Définition}

    \subsection*{Pseudo-code}
    \begin{verbatim}
Entrer epsilon
Entrer la limite supposée l
n <- 0
Tant que |u(n) - l| >= epsilon
    n <- n + 1
Fin Tant que
Afficher n
    \end{verbatim}

    \subsection*{Python (NumWorks ou PC)}
    \begin{verbatim}
def seuil(epsilon, l):
    n = 1
    u = 1/n   # exemple avec u_n = 1/n
    while abs(u - l) >= epsilon:
        n += 1
        u = 1/n
    return n

print(seuil(0.001, 0))  # renvoie 1000
    \end{verbatim}

    \chapter{Vecteurs, droites et plans de l'espace}
    \nchap

    \begin{capacitesbox}
    \begin{itemize}
        \item Représenter des droites et des plans dans l'espace et décrire leurs positions relatives.
        \item Utiliser les vecteurs pour caractériser le parallélisme et la coplanarité.
        \item Déterminer une représentation paramétrique d'une droite ou d'un plan.
        \item Se repérer dans un repère orthonormé de l'espace.
        \item Reconnaître, compléter ou construire une figure de l'espace.
    \end{itemize}
    \end{capacitesbox}

    \begin{demoprogbox}
    \begin{itemize}
        \item Si deux plans sécants contiennent chacun une droite de même direction, leur droite
              d'intersection est parallèle à ces droites.
        \item Une droite parallèle à deux droites sécantes d'un plan est parallèle à ce plan.
    \end{itemize}
    \end{demoprogbox}

    \section{Vecteurs de l'espace}
    \begin{Définition}
        Soient $A$ et $B$ deux points de l'espace. La transformation qui à tout point de
        l'espace associe l'unique point $M'$ tel que $ABM'M$ soit un parallélogramme est la
        \textbf{translation de vecteur} $\vc{AB}$.
        On note $\vc{u}=\vc{AB}=\vc{MM'}$ ; on dit que $\vc{AB}$ et $\vc{MM'}$ sont des
        \textbf{représentants} de $\vc{u}$.
    \end{Définition}

    \begin{Théorème}
        Soient un vecteur $\vc{u}$ et un point $A$ de l'espace. Il existe un unique point $M$
        tel que $\vc{AM}=\vc{u}$ ; on dit que $\vc{AM}$ est le représentant de $\vc{u}$
        d'origine $A$.
    \end{Théorème}

    \subsection{Opérations sur les vecteurs}
    \begin{Propriété}
        Soient $A$, $B$ et $C$ trois points. On a $\vc{AB}+\vc{AC}=\vc{AD}$ où $D$ est le
        point tel que $ABDC$ est un parallélogramme.
    \end{Propriété}

    \begin{Propriété}
        Soient $A$, $B$ et $C$ trois points. On a $\vc{AB}+\vc{BC}=\vc{AC}$ (relation de
        Chasles).
    \end{Propriété}

    \begin{Propriété}
        Soient $A$, $B$ des points distincts et $k\in\mathbf{R}$. Le vecteur $\vc{AC}=k\vc{AB}$
        est défini par :
        \begin{itemize}
            \item Si $k\geq 0$ alors $C\in[AB)$ et $AC=kAB$.
            \item Si $k\leq 0$ alors $C\in(AB)$ et $C\notin[AB)$, avec $AC=|k|AB$.
        \end{itemize}
    \end{Propriété}

    \begin{Propriété}
        Soient $\vc{u}$, $\vc{v}$ deux vecteurs et $k$, $k'$ deux réels.
        \begin{itemize}
            \item $k\vc{u}=\vc{0}$ si et seulement si $k=0$ ou $\vc{u}=\vc{0}$.
            \item $k(\vc{u}+\vc{v})=k\vc{u}+k\vc{v}$
            \item $(k+k')\vc{u}=k\vc{u}+k'\vc{u}$
        \end{itemize}
    \end{Propriété}

    \section{Droites et plans de l'espace}
    \begin{Définition}
        Soient $\vc{u}$ et $\vc{v}$ deux vecteurs de l'espace. On dit que $\vc{u}$ et $\vc{v}$
        sont \textbf{colinéaires} s'il existe un réel $k$ non nul tel que $\vc{u}=k\vc{v}$
        ou $\vc{v}=k\vc{u}$.
    \end{Définition}

    \begin{Propriété}
        Trois points $A$, $B$ et $C$ deux à deux distincts sont \textbf{alignés} si et
        seulement si $\vc{AB}$ et $\vc{AC}$ sont colinéaires.
    \end{Propriété}

    \begin{Propriété}
        \begin{enumerate}
            \item Par deux points distincts, il passe une unique droite.
            \item Par trois points non alignés $A$, $B$ et $C$, il passe un unique plan $(ABC)$.
            \item Si deux points distincts $A$ et $B$ appartiennent à un plan $\mathcal{P}$,
                  alors la droite $(AB)$ est incluse dans $\mathcal{P}$.
            \item Dans chaque plan de l'espace, toutes les règles de la géométrie plane
                  s'appliquent.
        \end{enumerate}
    \end{Propriété}

    \begin{Définition}
        Des vecteurs sont \textbf{coplanaires} si leurs représentants de même origine $A$ ont
        leurs extrémités dans un même plan passant par $A$.
    \end{Définition}

    \begin{Définition}
        Soient $A$ et $B$ deux points de l'espace.
        \begin{itemize}
            \item La droite $(AB)$ est l'ensemble des points $M$ tels que
                  $\vc{AM}=k\vc{AB}$ avec $k\in\mathbf{R}$.
            \item $\vc{AB}$ est un \textbf{vecteur directeur} de $(AB)$.
        \end{itemize}
    \end{Définition}

    \begin{Propriété}
        Une droite $(d)$ est définie par un point $A$ et un vecteur directeur $\vc{u}$ ;
        on la note $d(A;\vc{u})$.
    \end{Propriété}

    \begin{Définition}
        Soient $A$, $B$ et $C$ trois points de l'espace non alignés.
        \begin{itemize}
            \item Le plan $(ABC)$ est l'ensemble des points $M$ tels que
                  $\vc{AM}=x\vc{AB}+y\vc{AC}$ avec $x,y\in\mathbf{R}$.
            \item $\vc{AB}$ et $\vc{AC}$ sont des \textbf{vecteurs directeurs} du plan $(ABC)$.
        \end{itemize}
    \end{Définition}

    \begin{Propriété}
        Un plan $\mathcal{P}$ est défini par un point $A$ et deux vecteurs directeurs non
        colinéaires $\vc{u}$, $\vc{v}$ ; on le note $\mathcal{P}(A,\vc{u},\vc{v})$.
    \end{Propriété}

    \section{Positions relatives de droites et de plans}
    \subsection{Position de deux droites}
    \begin{Propriété}
        Deux droites $(d)$ et $(d')$ peuvent être coplanaires ou non coplanaires.
        \begin{center}
        \begin{tabular}{|c|c|c|}
            \hline
            \multicolumn{2}{|c|}{$(d)$ et $(d')$ sont coplanaires} & $(d)$ et $(d')$ sont non coplanaires \\
            \hline
            \psset{xunit=0.8cm,yunit=0.8cm}
            \begin{pspicture*}(-0.5,-0.5)(5,4.5)
                \pstGeonode[PointSymbol=x,PosAngle={-135,-45,-45,135,0,45,90,0}](1,1){A}(3,1){B}(3,3){F}(1,3){E}(4,1.5){C}(2,1.5){D}(2,3.5){H}(4,3.5){G}
                \pstLineAB{A}{B}\pstLineAB{F}{B}\pstLineAB{C}{B}\pstLineAB{A}{E}
                \pstLineAB{F}{G}\pstLineAB{C}{G}\pstLineAB{F}{E}\pstLineAB{H}{G}\pstLineAB{E}{H}
                \pstLineAB[linestyle=dashed]{A}{D}\pstLineAB[linestyle=dashed]{C}{D}\pstLineAB[linestyle=dashed]{D}{H}
                \pstLineAB[nodesepA=-2,nodesepB=-2]{E}{H}
                \pstLineAB[nodesepA=-2,nodesepB=-2]{F}{G}
            \end{pspicture*}
            &
            \begin{pspicture*}(-0.5,-0.5)(5,4.5)
                \psset{xunit=0.8cm,yunit=0.8cm}
                \pstGeonode[PointSymbol=x,PosAngle={-135,-45,-45,135,0,45,90,0}](1,1){A}(3,1){B}(3,3){F}(1,3){E}(4,1.5){C}(2,1.5){D}(2,3.5){H}(4,3.5){G}
                \pstLineAB{A}{B}\pstLineAB{F}{B}\pstLineAB{C}{B}\pstLineAB{A}{E}
                \pstLineAB{F}{G}\pstLineAB{C}{G}\pstLineAB{F}{E}\pstLineAB{H}{G}\pstLineAB{E}{H}
                \pstLineAB[linestyle=dashed]{A}{D}\pstLineAB[linestyle=dashed]{C}{D}\pstLineAB[linestyle=dashed]{D}{H}
                \pstLineAB[nodesepA=-2,nodesepB=-2]{E}{G}
                \pstLineAB[nodesepA=-2,nodesepB=-2]{F}{H}
            \end{pspicture*}
            &
            \begin{pspicture*}(-0.5,-0.5)(5,4.5)
                \psset{xunit=0.8cm,yunit=0.8cm}
                \pstGeonode[PointSymbol=x,PosAngle={-135,-45,-45,135,0,45,90,0}](1,1){A}(3,1){B}(3,3){F}(1,3){E}(4,1.5){C}(2,1.5){D}(2,3.5){H}(4,3.5){G}
                \pstLineAB{A}{B}\pstLineAB{F}{B}\pstLineAB{C}{B}\pstLineAB{A}{E}
                \pstLineAB{F}{G}\pstLineAB{C}{G}\pstLineAB{F}{E}\pstLineAB{H}{G}\pstLineAB{E}{H}
                \pstLineAB[linestyle=dashed]{A}{D}\pstLineAB[linestyle=dashed]{C}{D}\pstLineAB[linestyle=dashed]{D}{H}
                \pstLineAB[nodesepA=-2,nodesepB=-2]{E}{H}
                \pstLineAB[nodesepA=-2,nodesepB=-2]{G}{C}
            \end{pspicture*}
            \\
            $(EH)$ et $(FG)$ parallèles & $(EG)$ et $(FH)$ sécantes & $(EH)$ et $(GC)$ non coplanaires \\
            \hline
        \end{tabular}
        \end{center}
    \end{Propriété}

    \subsection{Position d'une droite et d'un plan}
    \begin{Propriété}
        Soit $(d)$ une droite et $\mathcal{P}$ un plan. Il existe trois positions relatives :
        \begin{center}
        \begin{tabular}{|c|c|c|}
            \hline
            $(d)$ strictement parallèle à $\mathcal{P}$ & $(d)$ incluse dans $\mathcal{P}$ & $(d)$ sécante à $\mathcal{P}$ \\
            \hline
            \begin{pspicture}(-0.5,-0.5)(5,4.5)
                \psset{xunit=0.8cm,yunit=0.8cm}
                \pstGeonode[PointSymbol=x,PosAngle={-135,-45,-45,135,0,45,90,0}](1,1){A}(3,1){B}(3,3){F}(1,3){E}(4,1.5){C}(2,1.5){D}(2,3.5){H}(4,3.5){G}
                \pstLineAB{A}{B}\pstLineAB{F}{B}\pstLineAB{C}{B}\pstLineAB{A}{E}
                \pstLineAB{F}{G}\pstLineAB{C}{G}\pstLineAB{F}{E}\pstLineAB{H}{G}\pstLineAB{E}{H}
                \pstLineAB[linestyle=dashed]{A}{D}\pstLineAB[linestyle=dashed]{C}{D}\pstLineAB[linestyle=dashed]{D}{H}
                \pstLineAB[nodesepA=-2,nodesepB=-2]{E}{H}
            \end{pspicture}
            &
            \begin{pspicture}(-0.5,-0.5)(5,4.5)
                \psset{xunit=0.8cm,yunit=0.8cm}
                \pstGeonode[PointSymbol=x,PosAngle={-135,-45,-45,135,0,45,90,0}](1,1){A}(3,1){B}(3,3){F}(1,3){E}(4,1.5){C}(2,1.5){D}(2,3.5){H}(4,3.5){G}
                \pstLineAB{A}{B}\pstLineAB{F}{B}\pstLineAB{C}{B}\pstLineAB{A}{E}
                \pstLineAB{F}{G}\pstLineAB{C}{G}\pstLineAB{F}{E}\pstLineAB{H}{G}\pstLineAB{E}{H}
                \pstLineAB[linestyle=dashed]{A}{D}\pstLineAB[linestyle=dashed]{C}{D}\pstLineAB[linestyle=dashed]{D}{H}
                \pstLineAB[nodesepA=-2,nodesepB=-2]{A}{C}
            \end{pspicture}
            &
            \begin{pspicture}(-0.5,-0.5)(5,4.5)
                \psset{xunit=0.8cm,yunit=0.8cm}
                \pstGeonode[PointSymbol=x,PosAngle={-135,-45,-45,135,0,45,90,0}](1,1){A}(3,1){B}(3,3){F}(1,3){E}(4,1.5){C}(2,1.5){D}(2,3.5){H}(4,3.5){G}
                \pstInterLL{E}{G}{H}{F}{O}
                \pstLineAB{A}{B}\pstLineAB{F}{B}\pstLineAB{C}{B}\pstLineAB{A}{E}
                \pstLineAB{F}{G}\pstLineAB{C}{G}\pstLineAB{F}{E}\pstLineAB{H}{G}\pstLineAB{E}{H}
                \pstLineAB[linestyle=dashed]{A}{D}\pstLineAB[linestyle=dashed]{C}{D}\pstLineAB[linestyle=dashed]{D}{H}
                \pstLineAB[nodesepA=-2,nodesepB=-2]{E}{H}
                \pstLineAB[nodesepA=-2,nodesepB=-2]{A}{O}
            \end{pspicture}
            \\
            $(EH) \parallel (ABC)$ & $(AC)$ incluse dans $(ABC)$ & $(AO)$ sécante à $(EFG)$ \\
            \hline
        \end{tabular}
        \end{center}
    \end{Propriété}

    \subsection{Position de deux plans}
    \begin{Propriété}
        Soient $\mathcal{P}$ et $\mathcal{P}'$ deux plans. Il existe trois configurations possibles :
        \begin{center}
        \begin{tabular}{|c|c|c|}
            \hline
            Plans parallèles & Plans sécants & Plans confondus \\
            \hline
            \begin{pspicture*}(-0.5,-0.5)(5,3)
                \psline(0,0)(2,0)(2,2)(0,2)(0,0)
                \psline(3,0.5)(3,2.5)(1,2.5)
                \psline(2,0.5)(3,0.5)
                \psline(1,2.5)(1,2)
                \psline[linestyle=dashed](1,2)(1,0.5)
                \psline[linestyle=dashed](2,0.5)(1,0.5)
            \end{pspicture*}
            &
            \begin{pspicture*}(-0.5,-1)(5,3)
                \psline(0,0)(2,0)(3,2)(1,2)(0,0)
                \psline(3.5,-0.5)(3,2)
                \psline[linestyle=dashed](1,2)(1.5,-0.5)
                \psline(1.5,-0.5)(3.5,-0.5)
            \end{pspicture*}
            &
            \begin{pspicture*}(-0.5,-0.5)(5,3)
                \psline(0,0)(2,0)(2,2)(0,2)(0,0)
                \psline(0.3,0.3)(2.3,0.3)(2.3,2.3)(0.3,2.3)(0.3,0.3)
            \end{pspicture*}
            \\
            \hline
        \end{tabular}
        \end{center}
    \end{Propriété}

    \begin{Théorème}
        Si deux plans $\mathcal{P}$ et $\mathcal{P}'$ contiennent respectivement deux droites
        $(d)$ et $(d')$ parallèles entre elles, alors leur intersection $(d'')$ est parallèle
        à $(d)$ et $(d')$.
    \end{Théorème}
    \begin{proof}
        Soit $\vc{v}$ un vecteur directeur commun à $(d)$ et $(d')$ (car elles sont parallèles).
        Soit $\vc{w}$ un vecteur directeur de $(d'')=({\mathcal{P}\cap\mathcal{P}'})$.
        Comme $(d'')$ est contenue dans $\mathcal{P}$, le vecteur $\vc{w}$ est une combinaison
        linéaire de $\vc{v}$ et d'un vecteur directeur de $(d)$.
        De même, $\vc{w}$ est une combinaison linéaire dans $\mathcal{P}'$.
        Si $\vc{w}$ et $\vc{v}$ n'étaient pas colinéaires, $(\vc{v},\vc{w})$ formerait une
        base commune à $\mathcal{P}$ et $\mathcal{P}'$, ce qui impliquerait $\mathcal{P}=\mathcal{P}'$,
        contradisant le fait qu'ils soient sécants. Donc $\vc{w}$ et $\vc{v}$ sont colinéaires,
        et $(d'')$ est parallèle à $(d)$ et $(d')$.
    \end{proof}

    \begin{Théorème}
        Si deux plans $\mathcal{P}$ et $\mathcal{P}'$ sont parallèles, alors tout plan $\Pi$
        qui coupe $\mathcal{P}$ coupe aussi $\mathcal{P}'$, et les intersections sont deux
        droites parallèles.
    \end{Théorème}

    \section{Repère de l'espace}
    \subsection{Base de l'espace}
    \begin{Définition}
        Trois vecteurs $\vc{i}$, $\vc{j}$ et $\vc{k}$ forment une \textbf{base de l'espace}
        si et seulement si aucun d'eux n'est combinaison linéaire des deux autres.
    \end{Définition}

    \begin{Propriété}
        Soit $(\vc{i},\vc{j},\vc{k})$ une base de l'espace. Tout vecteur $\vc{u}$ s'écrit
        de façon unique $\vc{u}=x\vc{i}+y\vc{j}+z\vc{k}$ ;
        $\left(\begin{smallmatrix}x\\y\\z\end{smallmatrix}\right)$ sont les coordonnées de
        $\vc{u}$ dans cette base.
    \end{Propriété}

    \subsection{Repère de l'espace}
    \begin{Propriété}
        Un \textbf{repère de l'espace} est un quadruplet $(O;\vc{i},\vc{j},\vc{k})$ où $O$
        est l'origine et $(\vc{i},\vc{j},\vc{k})$ est une base de l'espace.
    \end{Propriété}

    \begin{Propriété}
        Pour tout point $M$ de l'espace, il existe un unique triplet $(x_M;y_M;z_M)$ tel que
        $\vc{OM}=x_M\vc{i}+y_M\vc{j}+z_M\vc{k}$ ; ce triplet est appelé les
        \textbf{coordonnées} de $M$.
    \end{Propriété}

    \subsection{Coordonnées d'un vecteur}
    \begin{Propriété}
        Si $A(x_A;y_A;z_A)$ et $B(x_B;y_B;z_B)$ alors
        $\vc{AB}\vcb{x_B-x_A}{y_B-y_A}{z_B-z_A}$.
    \end{Propriété}

    \begin{Propriété}
        $\vc{u}=\vc{v}\Longleftrightarrow$ leurs coordonnées sont égales composante par composante.
    \end{Propriété}

    \begin{Propriété}
        Si $A(x_A;y_A;z_A)$ et $B(x_B;y_B;z_B)$, les coordonnées du milieu $I$ de $[AB]$ sont :
        \[
            x_I=\frac{x_A+x_B}{2},\quad y_I=\frac{y_A+y_B}{2},\quad z_I=\frac{z_A+z_B}{2}.
        \]
    \end{Propriété}

    \subsubsection{Opérations sur les coordonnées de vecteurs}
    \begin{Propriété}
        Soient $\vc{u}\vcb{x_u}{y_u}{z_u}$, $\vc{v}\vcb{x_v}{y_v}{z_v}$ et $\lambda\in\mathbf{R}$.
        \begin{enumerate}
            \item $\vc{u}+\vc{v}=\vcb{x_u+x_v}{y_u+y_v}{z_u+z_v}$
            \item $\lambda\vc{u}=\vcb{\lambda x_u}{\lambda y_u}{\lambda z_u}$
        \end{enumerate}
    \end{Propriété}

    \subsection{Représentation paramétrique d'une droite}
    \begin{Propriété}
        Dans un repère $\oijk$, la droite passant par $A(x_A;y_A;z_A)$ de vecteur directeur
        $\vc{u}\vcb{a}{b}{c}$ admet comme représentation paramétrique :
        \[
            \left\{\begin{array}{l}
                x=x_A+at\\
                y=y_A+bt\\
                z=z_A+ct
            \end{array}\right.\quad t\in\mathbf{R}.
        \]
    \end{Propriété}

    \chapter{Limites de fonctions}
    \nchap

    \begin{capacitesbox}
    \begin{itemize}
        \item Calculer des limites à l'aide des théorèmes d'opérations et de composition.
        \item Lever des formes indéterminées ($\infty-\infty$, $\frac{0}{0}$, $\frac{\infty}{\infty}$, $0\times\infty$).
        \item Utiliser les théorèmes de comparaison (gendarmes) et les croissances comparées.
        \item Étudier la continuité prolongeable en un point.
        \item Déterminer et interpréter les asymptotes horizontales, verticales et obliques.
    \end{itemize}
    \end{capacitesbox}

    \begin{demoprogbox}
    \begin{itemize}
        \item Démontrer $\displaystyle\lim_{x\to 0}\dfrac{\sin x}{x}=1$ à l'aide d'un encadrement géométrique.
        \item Démontrer l'unicité de la limite finie d'une fonction en un point (par l'absurde).
    \end{itemize}
    \end{demoprogbox}

    \section{Limite d'une fonction en l'infini}
    \begin{Définition}
        Une fonction $f$ a pour limite $+\infty$ en $+\infty$ si tout intervalle $]A;+\infty[$
        contient toutes les valeurs de $f(x)$ pour $x$ suffisamment grand.
        On note $\lim\limits_{x\to+\infty}f(x)=+\infty$.
    \end{Définition}

    \begin{Définition}
        Une fonction $f$ a pour limite $-\infty$ en $+\infty$ si tout intervalle $]-\infty;A[$
        contient toutes les valeurs de $f(x)$ pour $x$ suffisamment grand.
        On note $\lim\limits_{x\to+\infty}f(x)=-\infty$.
    \end{Définition}

    \begin{Définition}
        Une fonction $f$ a pour limite $\ell$ en $+\infty$ si tout intervalle ouvert contenant
        $\ell$ contient toutes les valeurs de $f(x)$ pour $x$ suffisamment grand.
        On note $\lim\limits_{x\to+\infty}f(x)=\ell$.
        La droite $y=\ell$ est une \textbf{asymptote horizontale} à $\mathcal{C}_f$ en $+\infty$.
    \end{Définition}

    \begin{Exemple}
        \begin{enumerate}
            \item $\lim\limits_{x\to+\infty}x^2=+\infty$
            \item $\lim\limits_{x\to+\infty}(-x^2)=-\infty$
            \item $\lim\limits_{x\to+\infty}\dfrac{2x+1}{x}=2$ car
                  $\dfrac{2x+1}{x}=2+\dfrac{1}{x}$ et $\lim\limits_{x\to+\infty}\dfrac{1}{x}=0$.
                  La droite $y=2$ est une asymptote horizontale.
        \end{enumerate}
    \end{Exemple}

    \section{Limite d'une fonction en un réel}
    \begin{Définition}
        Une fonction $f$ a pour limite $+\infty$ en $a$ si tout intervalle $]A;+\infty[$
        contient toutes les valeurs de $f(x)$ pour $x$ suffisamment proche de $a$.
        On note $\lim\limits_{x\to a}f(x)=+\infty$.
    \end{Définition}

    \begin{Définition}
        Une fonction $f$ a pour limite $\ell$ en $a$ si tout intervalle ouvert contenant
        $\ell$ contient toutes les valeurs de $f(x)$ pour $x$ suffisamment proche de $a$.
        On note $\lim\limits_{x\to a}f(x)=\ell$.
    \end{Définition}

    \begin{Définition}
        La droite $x=a$ est une \textbf{asymptote verticale} à $\mathcal{C}_f$ si au moins
        une des limites $\lim\limits_{x\to a^+}f(x)$ ou $\lim\limits_{x\to a^-}f(x)$ est
        infinie.
    \end{Définition}

    \begin{Définition}
        \begin{enumerate}
            \item On dit que $f$ admet une \textbf{limite à gauche} de $a$, notée
                  $\lim\limits_{x\to a^-}f(x)$, lorsque $f$ admet une limite quand $x\to a$
                  avec $x<a$.
            \item On dit que $f$ admet une \textbf{limite à droite} de $a$, notée
                  $\lim\limits_{x\to a^+}f(x)$, lorsque $f$ admet une limite quand $x\to a$
                  avec $x>a$.
        \end{enumerate}
    \end{Définition}

    \begin{Exemple}
        Pour $f(x)=\dfrac{1}{x-2}$ :
        \begin{enumerate}
            \item $\lim\limits_{x\to 2^+}f(x)=+\infty$
            \item $\lim\limits_{x\to 2^-}f(x)=-\infty$
            \item La droite $x=2$ est une asymptote verticale.
        \end{enumerate}
    \end{Exemple}

    \section{Opérations sur les limites}
    \subsection{Limite d'une somme}
    \begin{Propriété}
    \begin{center}
    \renewcommand{\arraystretch}{2}
    \begin{tabular}{|c|c|c|c|c|c|c|}
        \hline
        $\lim\limits_{x\to\alpha}f(x)$ & $\ell_1$ & $\ell_1$ & $\ell_1$ & $+\infty$ & $-\infty$ & $+\infty$ \\
        \hline
        $\lim\limits_{x\to\alpha}g(x)$ & $\ell_2$ & $+\infty$ & $-\infty$ & $+\infty$ & $-\infty$ & $-\infty$ \\
        \hline
        $\lim\limits_{x\to\alpha}(f+g)(x)$ & $\ell_1+\ell_2$ & $+\infty$ & $-\infty$ & $+\infty$ & $-\infty$ & \textbf{F.I.} \\
        \hline
    \end{tabular}
    \end{center}
    \end{Propriété}

    \subsection{Limite d'un produit}
    \begin{Propriété}
    \begin{center}
    \renewcommand{\arraystretch}{2}
    \begin{tabular}{|c|c|c|c|c|c|c|c|c|}
        \hline
        $\lim\limits_{x\to\alpha}f$ & $\ell_1$ & $\ell_1>0$ & $\ell_1>0$ & $\ell_1<0$ & $\ell_1<0$ & $+\infty$ & $-\infty$ & $0$ \\
        \hline
        $\lim\limits_{x\to\alpha}g$ & $\ell_2$ & $+\infty$ & $-\infty$ & $+\infty$ & $-\infty$ & $+\infty$ & $-\infty$ & $\pm\infty$ \\
        \hline
        $\lim\limits_{x\to\alpha}(f\times g)$ & $\ell_1\ell_2$ & $+\infty$ & $-\infty$ & $-\infty$ & $+\infty$ & $+\infty$ & $+\infty$ & \textbf{F.I.} \\
        \hline
    \end{tabular}
    \end{center}
    \end{Propriété}

    \subsection{Limite d'un quotient}
    \begin{Propriété}
    \begin{center}
    \renewcommand{\arraystretch}{2}
    \begin{tabular}{|c|c|c|c|c|c|c|}
        \hline
        $\lim\limits_{x\to\alpha}f$ & $\ell_1$ & $\ell_1$ & $\ell_1$ & $+\infty$ & $-\infty$ & $0$ \\
        \hline
        $\lim\limits_{x\to\alpha}g$ & $\ell_2\neq 0$ & $+\infty$ & $-\infty$ & $0^+$ & $0^-$ & $0$ \\
        \hline
        $\lim\limits_{x\to\alpha}\dfrac{f}{g}$ & $\dfrac{\ell_1}{\ell_2}$ & $0$ & $0$ & $+\infty$ & $-\infty$ & \textbf{F.I.} \\
        \hline
    \end{tabular}
    \end{center}
    \end{Propriété}

    \begin{Exemple}
        Calculons $\lim\limits_{x\to+\infty}\dfrac{x^2+1}{x^2-4x+1}$ :
        on divise numérateur et dénominateur par $x^2$ (terme de plus haut degré) :
        \[
            \lim\limits_{x\to+\infty}\frac{x^2+1}{x^2-4x+1}
            = \lim\limits_{x\to+\infty}\frac{1+\frac{1}{x^2}}{1-\frac{4}{x}+\frac{1}{x^2}}
            = \frac{1+0}{1-0+0} = 1.
        \]
    \end{Exemple}

    \section{Composition de limites}
    \begin{Théorème}
        Si $\lim\limits_{x\to a}f(x)=b$ et $\lim\limits_{x\to b}g(x)=c$, alors
        $\lim\limits_{x\to a}g(f(x))=c$.
    \end{Théorème}

    \begin{Exemple}
        Soit $f(x)=x^2$ et $g(x)=\sqrt{x}$.
        $\lim\limits_{x\to 2}f(x)=4$ et $\lim\limits_{x\to 4}g(x)=2$, donc
        $\lim\limits_{x\to 2}g(f(x))=\lim\limits_{x\to 2}\sqrt{x^2}=2$.
    \end{Exemple}

    \section{Théorèmes de comparaison}
    \begin{Théorème}[Théorème de comparaison]
        \begin{enumerate}
            \item Si pour $x$ suffisamment grand $f(x)\leq g(x)$ et
                  $\lim\limits_{x\to+\infty}f(x)=+\infty$, alors
                  $\lim\limits_{x\to+\infty}g(x)=+\infty$.
            \item Si pour $x$ suffisamment grand $f(x)\leq g(x)$ et
                  $\lim\limits_{x\to+\infty}g(x)=-\infty$, alors
                  $\lim\limits_{x\to+\infty}f(x)=-\infty$.
        \end{enumerate}
    \end{Théorème}

    \begin{Théorème}[Théorème d'encadrement]
        Si pour $x$ suffisamment grand $f(x)\leq g(x)\leq h(x)$ et
        $\lim\limits_{x\to+\infty}f(x)=\lim\limits_{x\to+\infty}h(x)=\ell$, alors
        $\lim\limits_{x\to+\infty}g(x)=\ell$.
    \end{Théorème}

    \begin{Théorème}[Croissances comparées]
        Pour tout $n\in\mathbb{N}$ :
        $\lim\limits_{x\to+\infty}\dfrac{e^x}{x^n}=+\infty$ et
        $\lim\limits_{x\to-\infty}x^n e^x=0$.
    \end{Théorème}

    \chapter{Dérivation et convexité}
    \nchap

    \begin{capacitesbox}
    \begin{itemize}
        \item Calculer la dérivée d'une fonction composée $v\circ u$ : $(v\circ u)'(x)=u'(x)\,v'(u(x))$.
        \item Caractériser la convexité d'une fonction par le signe de $f''$.
        \item Exploiter la convexité pour comparer $f(x)$ avec $f(a)+f'(a)(x-a)$ (tangente).
        \item Déterminer les points d'inflexion d'une courbe.
        \item Appliquer la convexité pour établir des inégalités classiques.
    \end{itemize}
    \end{capacitesbox}

    \begin{demoprogbox}
    \begin{itemize}
        \item Démontrer la règle de dérivation en chaîne : $[v(u(x))]'=u'(x)\,v'(u(x))$.
        \item Démontrer que si $f''\geq 0$ sur $I$, la courbe est au-dessus de chacune de ses tangentes.
    \end{itemize}
    \end{demoprogbox}

    \section{Compléments sur les dérivées}
    \subsection{Dérivée d'une fonction composée}
    \begin{Définition}
        Soient $u$ définie sur $I$ à valeurs dans $J$ et $v$ définie sur $J$.
        La \textbf{composée} $v\circ u$ est définie sur $I$ par $(v\circ u)(x)=v(u(x))$.
    \end{Définition}

    \begin{Propriété}
        Si $u$ est dérivable sur $I$ et $v$ est dérivable sur $J$, alors $v\circ u$ est
        dérivable sur $I$ et :
        \[
            (v\circ u)'(x_0)=u'(x_0)\times v'(u(x_0)).
        \]
    \end{Propriété}

    \begin{Propriété}
        \begin{enumerate}
            \item Si $u$ est dérivable sur $I$, la fonction $f(x)=e^{u(x)}$ est dérivable
                  et $f'(x)=u'(x)\,e^{u(x)}$.
            \item Si $u$ est dérivable sur $I$ à valeurs dans $]0;+\infty[$, la fonction
                  $f(x)=\sqrt{u(x)}$ est dérivable et $f'(x)=\dfrac{u'(x)}{2\sqrt{u(x)}}$.
        \end{enumerate}
    \end{Propriété}

    \begin{Remarque}
        En général $u\circ v\neq v\circ u$. Par exemple avec $u(x)=x^2$ et $v(x)=x+1$ :
        $u(v(x))=(x+1)^2=x^2+2x+1$ mais $v(u(x))=x^2+1$.
    \end{Remarque}

    \subsection{Dérivée seconde}
    \begin{Définition}
        Soit $f$ dérivable sur $I$. Si $f'$ est elle-même dérivable sur $I$, sa dérivée est
        notée $f''$ et appelée \textbf{dérivée seconde} de $f$.
    \end{Définition}

    \section{Convexité}
    \subsection{Fonctions convexes et concaves}
    \begin{Définition}
        Soit $f$ définie sur $I$ et $\mathcal{C}_f$ sa courbe représentative. Soient $A$ et
        $B$ deux points de $\mathcal{C}_f$ ; la droite $(AB)$ est appelée \textbf{sécante}
        de $\mathcal{C}_f$.
    \end{Définition}

    \begin{Définition}
        \begin{enumerate}
            \item $f$ est \textbf{convexe} sur $I$ si $\mathcal{C}_f$ est en dessous de
                  chacune de ses sécantes entre les deux points d'intersection.
            \item $f$ est \textbf{concave} sur $I$ si $\mathcal{C}_f$ est au-dessus de
                  chacune de ses sécantes entre les deux points d'intersection.
        \end{enumerate}
    \end{Définition}

    \begin{Propriété}
        Si $f$ est deux fois dérivable sur $I$, les propositions suivantes sont équivalentes :
        \begin{enumerate}
            \item $f$ est convexe sur $I$ ;
            \item $f'$ est croissante sur $I$ ;
            \item $f''(x)\geq 0$ pour tout $x\in I$.
        \end{enumerate}
    \end{Propriété}

    \begin{Propriété}
        Si $f''(x)\geq 0$ sur $I$, alors $\mathcal{C}_f$ est au-dessus de chacune de ses
        tangentes.
    \end{Propriété}

    \subsection{Point d'inflexion}
    \begin{Définition}
        Un \textbf{point d'inflexion} est un point où la courbe traverse sa tangente.
    \end{Définition}

    \begin{Propriété}
        Le point $A(a;f(a))$ est un point d'inflexion si et seulement si $f''$ s'annule en
        $a$ en changeant de signe.
    \end{Propriété}

    \chapter{Continuité}
    \nchap

    \begin{capacitesbox}
    \begin{itemize}
        \item Comprendre et utiliser la définition de la continuité en un point et sur un intervalle.
        \item Reconnaître et utiliser la continuité des fonctions usuelles et de leurs composées.
        \item Appliquer le théorème des valeurs intermédiaires et son corollaire (monotonie stricte).
        \item Encadrer une solution de $f(x)=k$ par la méthode de dichotomie.
        \item Étudier la convergence d'une suite récurrente $u_{n+1}=f(u_n)$ via le point fixe.
    \end{itemize}
    \end{capacitesbox}

    \begin{demoprogbox}
    \begin{itemize}
        \item Démontrer un cas particulier du TVI : si $f$ est continue sur $[a;b]$, $f(a)<0$
              et $f(b)>0$, alors $f$ s'annule sur $]a;b[$ (admis en général, cas particulier exigible).
        \item Démontrer que la méthode de dichotomie fournit une suite encadrant la solution avec
              une précision $\dfrac{b-a}{2^n}$.
    \end{itemize}
    \end{demoprogbox}

    \section{Fonctions continues sur un intervalle}
    \begin{Définition}
        Soit $f$ définie sur un intervalle $I$ et $a\in I$.
        \begin{itemize}
            \item $f$ est \textbf{continue en $a$} si $\lim\limits_{x\to a}f(x)=f(a)$.
            \item $f$ est \textbf{continue sur $I$} si elle est continue en tout réel $a\in I$.
        \end{itemize}
    \end{Définition}

    \section{Continuité de fonctions usuelles}
    \begin{Propriété}
    \begin{itemize}
        \item Les fonctions affines et polynômes sont continues sur $\mathbf{R}$.
        \item La fonction exponentielle est continue sur $\mathbf{R}$.
        \item La fonction racine carrée est continue sur $[0;+\infty[$.
    \end{itemize}
    \end{Propriété}

    \begin{Propriété}
        Les sommes, produits, quotients et composées de fonctions continues sont continues
        sur leur domaine de définition.
    \end{Propriété}

    \begin{Propriété}
        Toute fonction dérivable sur $I$ est continue sur $I$.
    \end{Propriété}

    \begin{Remarque}
        La réciproque est fausse : la fonction valeur absolue est continue en $0$ mais n'y
        est pas dérivable.
    \end{Remarque}

    \section{Théorème des valeurs intermédiaires}
    \begin{Théorème}
        Si $f$ est continue sur $[a;b]$, alors pour tout réel $k$ compris entre $f(a)$ et
        $f(b)$, il existe au moins un réel $c\in[a;b]$ tel que $f(c)=k$.
    \end{Théorème}

    \begin{Théorème}[Corollaire]
        Si $f$ est continue et strictement monotone sur $[a;b]$, alors pour tout $k$ compris
        entre $f(a)$ et $f(b)$, l'équation $f(x)=k$ possède une \textbf{unique} solution
        dans $[a;b]$.
    \end{Théorème}

    \section{Fonctions continues et suites convergentes}
    \begin{Propriété}
        Si $(u_n)$ est une suite à valeurs dans $I$, convergente vers $\ell\in I$, et $f$ est
        continue sur $I$, alors $\lim\limits_{n\to+\infty}f(u_n)=f(\ell)$.
    \end{Propriété}

    \begin{Théorème}[Théorème du point fixe]
        Soient $f$ continue sur $I$ dans lui-même et $(u_n)$ définie par $u_0\in I$ et
        $u_{n+1}=f(u_n)$.
        Si $(u_n)$ converge vers $\ell\in I$, alors $\ell$ est solution de $f(x)=x$.
    \end{Théorème}

    \chapter{Orthogonalité, équations et distances dans l'espace}
    \nchap

    \begin{capacitesbox}
    \begin{itemize}
        \item Calculer un produit scalaire dans l'espace à partir des coordonnées ou des normes.
        \item Reconnaître ou démontrer l'orthogonalité de deux vecteurs, deux droites, une droite et un plan.
        \item Déterminer une équation cartésienne d'un plan à partir d'un vecteur normal.
        \item Calculer la distance d'un point à un plan.
        \item Utiliser le produit scalaire pour calculer des longueurs et des angles dans l'espace.
    \end{itemize}
    \end{capacitesbox}

    \begin{demoprogbox}
    \begin{itemize}
        \item Démontrer l'expression du produit scalaire en coordonnées :
              $\vc{u}\cdot\vc{v}=xx'+yy'+zz'$.
        \item Démontrer la formule de la distance d'un point $A(x_0,y_0,z_0)$ au plan d'équation
              $ax+by+cz+d=0$ : $d=\dfrac{|ax_0+by_0+cz_0+d|}{\sqrt{a^2+b^2+c^2}}$.
    \end{itemize}
    \end{demoprogbox}

    \section{Produit scalaire dans l'espace}
    \begin{Définition}
        Soient $\vc{u}=\vc{AB}$ et $\vc{v}=\vc{AC}$ deux vecteurs de l'espace.
        Les points $A$, $B$, $C$ déterminent un plan $\mathcal{P}$.
        Le \textbf{produit scalaire} $\vc{u}\cdot\vc{v}$ est le produit scalaire de
        $\vc{AB}$ et $\vc{AC}$ dans ce plan.
    \end{Définition}

    \begin{Propriété}
        Soient $\vc{u}=\vc{AB}$ et $\vc{v}=\vc{AC}$ deux vecteurs non nuls.
        \begin{enumerate}
            \item $\vc{u}\cdot\vc{v}=AB\times AC\times\cos(\widehat{BAC})$
            \item Si $H$ est le projeté orthogonal de $C$ sur $(AB)$ alors
                  $\vc{u}\cdot\vc{v}=\vc{AB}\cdot\vc{AH}$.
        \end{enumerate}
    \end{Propriété}

    \begin{Définition}
        Deux vecteurs $\vc{u}$ et $\vc{v}$ sont \textbf{orthogonaux} si $\vc{u}\cdot\vc{v}=0$.
    \end{Définition}

    \subsection{Opérations avec le produit scalaire}
    \begin{Propriété}
        \begin{enumerate}
            \item $\vc{u}\cdot\vc{v}=\vc{v}\cdot\vc{u}$
            \item $\vc{u}\cdot(k\vc{v})=k(\vc{u}\cdot\vc{v})$
            \item $\vc{u}\cdot(\vc{v}+\vc{w})=\vc{u}\cdot\vc{v}+\vc{u}\cdot\vc{w}$
            \item $\|\vc{u}+\vc{v}\|^2=\|\vc{u}\|^2+2\,\vc{u}\cdot\vc{v}+\|\vc{v}\|^2$
            \item $\|\vc{u}-\vc{v}\|^2=\|\vc{u}\|^2-2\,\vc{u}\cdot\vc{v}+\|\vc{v}\|^2$
            \item $\vc{u}\cdot\vc{v}=\dfrac{1}{2}\!\left(\|\vc{u}\|^2+\|\vc{v}\|^2-\|\vc{u}-\vc{v}\|^2\right)$
            \item $\vc{u}\cdot\vc{v}=\dfrac{1}{2}\!\left(\|\vc{u}+\vc{v}\|^2-\|\vc{u}\|^2-\|\vc{v}\|^2\right)$
        \end{enumerate}
    \end{Propriété}

    \section{Produit scalaire dans un repère}
    \begin{Définition}
        Une base $(\vc{i},\vc{j},\vc{k})$ est \textbf{orthonormée} si les vecteurs sont deux
        à deux orthogonaux et de norme $1$.
        Un repère $(O;\vc{i},\vc{j},\vc{k})$ est \textbf{orthonormé} si sa base l'est.
    \end{Définition}

    \begin{Propriété}
        Dans un repère orthonormé $\oijk$, soient $\vc{u}\vcb{x}{y}{z}$ et
        $\vc{v}\vcb{x'}{y'}{z'}$.
        \begin{enumerate}
            \item $\vc{u}\cdot\vc{v}=xx'+yy'+zz'$
            \item $\|\vc{u}\|=\sqrt{x^2+y^2+z^2}$
            \item $AB=\sqrt{(x_B-x_A)^2+(y_B-y_A)^2+(z_B-z_A)^2}$
        \end{enumerate}
    \end{Propriété}

    \section{Orthogonalité dans l'espace}
    \subsection{Orthogonalité de deux droites}
    \begin{Définition}
        Deux droites sont \textbf{orthogonales} si leurs vecteurs directeurs sont orthogonaux.
    \end{Définition}

    \begin{Propriété}
        $(d_1)$ de vecteur directeur $\vc{u}$ et $(d_2)$ de vecteur directeur $\vc{v}$ sont
        orthogonales si et seulement si $\vc{u}\cdot\vc{v}=0$.
    \end{Propriété}

    \subsection{Orthogonalité d'une droite et d'un plan}
    \begin{Définition}
        Une droite est \textbf{orthogonale} à un plan si elle est orthogonale à toute droite
        de ce plan.
    \end{Définition}

    \begin{Propriété}
        $(d)$ de vecteur directeur $\vc{w}$ est orthogonale au plan $\mathcal{P}$ dirigé par
        $(\vc{u},\vc{v})$ si et seulement si $\vc{u}\cdot\vc{w}=0$ et $\vc{v}\cdot\vc{w}=0$.
    \end{Propriété}

    \subsection{Équation cartésienne d'un plan}
    \begin{Définition}
        Un \textbf{vecteur normal} au plan $\mathcal{P}$ est tout vecteur non nul orthogonal
        à tous les vecteurs directeurs de $\mathcal{P}$.
    \end{Définition}

    \begin{Propriété}
        L'ensemble des points $M$ tels que $\vc{AM}\cdot\vc{n}=0$ est le plan passant par
        $A$ de vecteur normal $\vc{n}$.
    \end{Propriété}

    \begin{Propriété}
        Dans un repère orthonormé, le plan passant par $A(x_A;y_A;z_A)$ de vecteur normal
        $\vc{n}\vcb{a}{b}{c}$ a pour équation cartésienne :
        \[
            ax+by+cz+d=0 \quad\text{avec}\quad d=-(ax_A+by_A+cz_A).
        \]
    \end{Propriété}

    \section{Projection orthogonale et distances}
    \begin{Définition}
        Le \textbf{projeté orthogonal} d'un point $A$ sur une droite (ou un plan) est le pied
        de la perpendiculaire menée de $A$ à cette droite (ou ce plan).
    \end{Définition}

    \begin{Définition}
        La \textbf{distance} d'un point $A$ à un plan $\mathcal{P}$ est la plus petite
        longueur $AM$ pour $M\in\mathcal{P}$.
    \end{Définition}

    \begin{Propriété}
        Si $H$ est le projeté orthogonal de $A$ sur $\mathcal{P}$ et $B$ un point de
        $\mathcal{P}$, alors $d(A,\mathcal{P})=AH$ et :
        \[
            AH=\frac{|\vc{AB}\cdot\vc{n}|}{\|\vc{n}\|}.
        \]
    \end{Propriété}

    \begin{Propriété}
        Si $H$ est le projeté orthogonal de $A$ sur la droite $(d)$ de vecteur directeur
        $\vc{u}$ et $B$ un point de $(d)$, alors $d(A,d)=AH$ et :
        \[
            AH=\left\|\vc{AB}-\frac{\vc{AB}\cdot\vc{u}}{\|\vc{u}\|^2}\,\vc{u}\right\|.
        \]
    \end{Propriété}

    \chapter{Loi binomiale}
    \nchap

    \begin{capacitesbox}
    \begin{itemize}
        \item Identifier une expérience de Bernoulli et une loi binomiale $\mathcal{B}(n;p)$.
        \item Calculer des probabilités $P(X=k)$, $P(X\leq k)$, $P(X\geq k)$ pour $X\sim\mathcal{B}(n;p)$.
        \item Calculer l'espérance, la variance et l'écart-type de $X\sim\mathcal{B}(n;p)$.
        \item Utiliser la loi binomiale pour modéliser une situation probabiliste.
    \end{itemize}
    \end{capacitesbox}

    \begin{demoprogbox}
    \begin{itemize}
        \item Démontrer que le nombre de parties d'un ensemble à $n$ éléments est $2^n$.
        \item Démontrer la formule $P(X=k)=\dbinom{n}{k}p^k(1-p)^{n-k}$ (chemins favorables dans un arbre).
    \end{itemize}
    \end{demoprogbox}

    \section{Dénombrement}
    \subsection{Principe additif}
    \begin{Définition}
        Deux ensembles $A$ et $B$ sont \textbf{disjoints} si $A\cap B=\emptyset$.
    \end{Définition}

    \begin{Définition}
        Un ensemble $E$ à $n$ éléments est dit \textbf{fini} de \textbf{cardinal} $n$,
        noté $\mathrm{Card}(E)=n$.
    \end{Définition}

    \begin{Propriété}
        Si $A_1,\ldots,A_p$ sont des ensembles deux à deux disjoints :
        \[
            \mathrm{Card}(A_1\cup\cdots\cup A_p)=\mathrm{Card}(A_1)+\cdots+\mathrm{Card}(A_p).
        \]
    \end{Propriété}

    \subsection{Principe multiplicatif}
    \begin{Définition}
        Un \textbf{$p$-uplet} de $E_1\times\cdots\times E_p$ est une liste ordonnée
        $(x_1;\ldots;x_p)$ avec $x_i\in E_i$. Un 2-uplet est un couple, un 3-uplet un triplet.
    \end{Définition}

    \begin{Propriété}
        $\mathrm{Card}(E_1\times\cdots\times E_p)=\mathrm{Card}(E_1)\times\cdots\times
        \mathrm{Card}(E_p)$.
    \end{Propriété}

    \section{Arrangements et permutations}
    \begin{Définition}
        On appelle \textbf{factorielle} de $n$ : $n!=n\times(n-1)\times\cdots\times 1$.
        Par convention $0!=1$.
    \end{Définition}

    \begin{Définition}
        Un \textbf{arrangement} de $k$ éléments d'un ensemble $A$ à $n$ éléments est un
        $k$-uplet d'éléments \emph{distincts} de $A$.
    \end{Définition}

    \begin{Propriété}
        Le nombre d'arrangements de $k$ éléments parmi $n$ est :
        \[
            n\times(n-1)\times\cdots\times(n-k+1)=\frac{n!}{(n-k)!}.
        \]
    \end{Propriété}

    \begin{Définition}
        Une \textbf{permutation} d'un ensemble $E$ à $n$ éléments est un $n$-uplet
        d'éléments distincts de $E$.
    \end{Définition}

    \begin{Propriété}
        Le nombre de permutations d'un ensemble à $n$ éléments est $n!$.
    \end{Propriété}

    \section{Combinaisons}
    \begin{Propriété}
        Le nombre de parties d'un ensemble fini à $n$ éléments est $2^n$.
    \end{Propriété}
    \begin{proof}
        Chaque élément peut être inclus (1) ou exclu (0) d'une partie donnée.
        Il y a donc $2^n$ choix, correspondant aux $2^n$ suites binaires de longueur $n$.
    \end{proof}

    \subsection{Nombre de combinaisons}
    \begin{Définition}
        Une \textbf{combinaison} de $k$ éléments de $E$ est une partie de $E$ à $k$ éléments.
        Le nombre de telles combinaisons est noté $\dbinom{n}{k}$.
    \end{Définition}

    \begin{Propriété}
        Pour $k\leq n$ :
        \begin{enumerate}
            \item $\dbinom{n}{k}=\dfrac{n!}{k!\,(n-k)!}$ et $\dbinom{n}{k}=\dbinom{n}{n-k}$.
            \item \textbf{Relation de Pascal :} si $1\leq k\leq n-1$,
                  $\dbinom{n}{k}=\dbinom{n-1}{k-1}+\dbinom{n-1}{k}$.
            \item $\dbinom{n}{0}=1$, $\dbinom{n}{1}=n$, $\dbinom{n}{2}=\dfrac{n(n-1)}{2}$.
        \end{enumerate}
    \end{Propriété}

    \section{Successions d'épreuves indépendantes}
    \begin{Définition}
        L'univers d'une succession de $n$ épreuves indépendantes d'univers
        $\Omega_1,\ldots,\Omega_n$ est $\Omega_1\times\cdots\times\Omega_n$.
    \end{Définition}

    \begin{Propriété}
        La probabilité d'une issue $(x_1;\ldots;x_n)$ est
        $p(x_1)\times\cdots\times p(x_n)$.
    \end{Propriété}

    \section{Loi et schéma de Bernoulli}
    \begin{Définition}
        Une \textbf{épreuve de Bernoulli} est une expérience à deux issues : succès $S$ et
        échec $\bar{S}$.
    \end{Définition}

    \begin{Définition}
        La \textbf{loi de Bernoulli} de paramètre $p$ est la loi de la variable $X$ :
        \begin{center}
        \begin{tabular}{|c|c|c|}
        \hline
        $x_i$ & $0$ & $1$ \\
        \hline
        $P(X=x_i)$ & $1-p$ & $p$ \\
        \hline
        \end{tabular}
        \end{center}
        On note $X\sim\mathcal{B}(p)$.
    \end{Définition}

    \begin{Propriété}
        Pour $X\sim\mathcal{B}(p)$ : $E(X)=p$, $V(X)=p(1-p)$,
        $\sigma(X)=\sqrt{p(1-p)}$.
    \end{Propriété}

    \begin{Définition}
        Un \textbf{schéma de Bernoulli} est la répétition de $n$ épreuves de Bernoulli
        identiques et indépendantes.
    \end{Définition}

    \begin{Définition}
        La variable $X$ comptant le nombre de succès dans un schéma de Bernoulli de
        paramètres $n$ et $p$ suit la \textbf{loi binomiale} $\mathcal{B}(n;p)$.
    \end{Définition}

    \begin{Propriété}
        Pour $X\sim\mathcal{B}(n;p)$ et $k\in\{0,\ldots,n\}$ :
        \[
            P(X=k)=\binom{n}{k}p^k(1-p)^{n-k}.
        \]
    \end{Propriété}
    \begin{proof}
        Un chemin réalisant $k$ succès et $n-k$ échecs a pour probabilité $p^k(1-p)^{n-k}$.
        Le nombre de tels chemins est $\dbinom{n}{k}$ (nombre de façons de placer $k$ succès
        parmi $n$ épreuves). Donc $P(X=k)=\dbinom{n}{k}p^k(1-p)^{n-k}$.
    \end{proof}

    \begin{Propriété}
        Pour $X\sim\mathcal{B}(n;p)$ : $E(X)=np$, $V(X)=np(1-p)$,
        $\sigma(X)=\sqrt{np(1-p)}$.
    \end{Propriété}

    \chapter{Fonction logarithme népérien}
    \nchap

    \begin{capacitesbox}
    \begin{itemize}
        \item Utiliser les propriétés algébriques de $\ln$ : $\ln(ab)=\ln a+\ln b$, $\ln(a/b)=\ln a-\ln b$, $\ln(a^n)=n\ln a$.
        \item Résoudre des équations et inéquations faisant intervenir $\ln$.
        \item Étudier (variations, limites, asymptotes) des fonctions contenant $\ln$.
        \item Utiliser les croissances comparées : $\ln x\ll x^\alpha$ pour tout $\alpha>0$.
    \end{itemize}
    \end{capacitesbox}

    \begin{demoprogbox}
    \begin{itemize}
        \item Démontrer que $(\ln x)'=\dfrac{1}{x}$ (par la règle de dérivation des fonctions réciproques).
        \item Démontrer $\displaystyle\lim_{x\to+\infty}\dfrac{\ln x}{x}=0$ (croissances comparées).
        \item Démontrer $\displaystyle\lim_{x\to 0^+}x\ln x=0$.
    \end{itemize}
    \end{demoprogbox}

    \begin{Définition}
        La \textbf{fonction logarithme népérien}, notée $\ln$, est définie sur $]0;+\infty[$
        et associe à tout $x>0$ l'unique solution $y$ de l'équation $e^y=x$.
    \end{Définition}

    \begin{Propriété}
        Les courbes des fonctions exponentielle et $\ln$ sont symétriques par rapport à la
        droite $y=x$.
    \end{Propriété}

    \section{Propriétés algébriques}
    \begin{Propriété}
        Pour tout $a>0$, $b>0$ et $x\in\mathbf{R}$ :
        \begin{enumerate}
            \item $\ln(e^x)=x$
            \item $e^{\ln x}=x$
            \item $\ln(ab)=\ln a+\ln b$
            \item $\ln\!\left(\dfrac{a}{b}\right)=\ln a-\ln b$
            \item $\ln(a^n)=n\ln a$ pour $n\in\mathbb{Z}$
            \item $\ln(\sqrt{a})=\dfrac{1}{2}\ln a$
            \item $\ln\!\left(\dfrac{1}{x}\right)=-\ln x$
        \end{enumerate}
    \end{Propriété}

    \begin{Remarque}
        Valeurs particulières : $\ln 1=0$ et $\ln e=1$.
    \end{Remarque}

    \section{Équations avec logarithme et exponentielle}
    \subsection{Équation de la forme $\ln(X)=a$}
    La solution unique est $X=e^a$.

    \subsection{Équation de la forme $e^X=a$}
    \begin{enumerate}
        \item Si $a\leq 0$, il n'y a pas de solution.
        \item Si $a>0$, la solution est $X=\ln a$.
    \end{enumerate}

    \section{Étude de la fonction logarithme}
    \subsection{Continuité et limites}
    \begin{Propriété}
        La fonction $\ln$ est continue sur $]0;+\infty[$.
    \end{Propriété}

    \begin{Propriété}
        \textbf{Limites aux bornes :}
        $\lim\limits_{x\to 0^+}\ln x=-\infty$ et
        $\lim\limits_{x\to+\infty}\ln x=+\infty$.\\[4pt]
        \textbf{Croissances comparées :} pour tout entier $k\geq 1$,
        \[
            \lim\limits_{x\to+\infty}\frac{\ln x}{x^k}=0
            \qquad\text{et}\qquad
            \lim\limits_{x\to 0^+}x^k\ln x=0.
        \]
    \end{Propriété}

    \subsection{Dérivée de la fonction logarithme}
    \begin{Propriété}
        La fonction $\ln$ est dérivable sur $]0;+\infty[$ avec $\ln'(x)=\dfrac{1}{x}$.\\
        Si $u$ est dérivable sur $I$ et $u(x)>0$ pour tout $x\in I$, alors
        $(\ln\circ\, u)'(x)=\dfrac{u'(x)}{u(x)}$.
    \end{Propriété}

    \chapter{Primitives et équations différentielles}
    \nchap

    \begin{capacitesbox}
    \begin{itemize}
        \item Calculer des primitives des fonctions usuelles (tableau de primitives).
        \item Résoudre les équations différentielles $y'=ay$ et $y'=ay+b$ ($a\neq 0$).
        \item Résoudre un problème de Cauchy (condition initiale).
        \item Modéliser un phénomène physique ou biologique par une équation différentielle.
    \end{itemize}
    \end{capacitesbox}

    \begin{demoprogbox}
    \begin{itemize}
        \item Démontrer que les fonctions $x\mapsto Ce^{ax}$ sont les seules solutions de $y'=ay$.
        \item Démontrer que les solutions de $y'=ay+b$ sont les fonctions $x\mapsto Ce^{ax}-\dfrac{b}{a}$.
    \end{itemize}
    \end{demoprogbox}

    \begin{Définition}
    \begin{itemize}
        \item Une \textbf{équation différentielle} est une égalité liant une fonction
              inconnue $y$ et ses dérivées successives.
        \item Une \textbf{solution} est toute fonction dérivable vérifiant cette égalité.
    \end{itemize}
    \end{Définition}

    \begin{Définition}
        $F$ est une \textbf{primitive} de $f$ sur $I$ si $F$ est dérivable sur $I$ et
        $F'=f$.
    \end{Définition}

    \subsection{Primitives de fonctions usuelles}
    \renewcommand{\arraystretch}{2}
    \begin{tabular}{|m{5cm}|m{4cm}|m{7cm}|}
        \hline
        Fonction $f$ & Domaine & Primitive $F$ \\
        \hline
        $f(x)=a$ & $\mathbf{R}$ & $F(x)=ax+c$ \\
        \hline
        $f(x)=x^n$, $n\in\mathbb{Z}\setminus\{-1\}$ &
            $\mathbf{R}$ si $n\in\mathbf{N}$;\; $\mathbf{R}^*$ sinon &
            $F(x)=\dfrac{x^{n+1}}{n+1}+c$ \\
        \hline
        $f(x)=\dfrac{1}{\sqrt{x}}$ & $]0;+\infty[$ & $F(x)=2\sqrt{x}+c$ \\
        \hline
        $f(x)=e^x$ & $\mathbf{R}$ & $F(x)=e^x+c$ \\
        \hline
        $f(x)=\dfrac{1}{x}$ & $]0;+\infty[$ ou $]-\infty;0[$ & $F(x)=\ln|x|+c$ \\
        \hline
    \end{tabular}

    \begin{Théorème}
        Toute fonction continue sur un intervalle $I$ admet une primitive sur $I$.
    \end{Théorème}

    \begin{Propriété}
        Deux primitives d'une même fonction continue sur $I$ diffèrent d'une constante.
    \end{Propriété}

    \begin{Propriété}
        Pour tout $(x_0,y_0)\in I\times\mathbf{R}$, il existe une unique solution $F$ de
        $y'=f$ sur $I$ vérifiant $F(x_0)=y_0$.
    \end{Propriété}

    \begin{Propriété}
        Si $F$ est une primitive de $f$ et $G$ de $g$ sur $I$, et $\lambda\in\mathbf{R}$ :
        \begin{itemize}
            \item $F+G$ est une primitive de $f+g$ sur $I$.
            \item $\lambda F$ est une primitive de $\lambda f$ sur $I$.
        \end{itemize}
    \end{Propriété}

    \section{Primitives de fonctions composées}
    \renewcommand{\arraystretch}{2.5}
    \begin{tabular}{|m{5cm}|m{4cm}|m{7cm}|}
        \hline
        Forme de $f$ & Primitive (sans constante) & Conditions \\
        \hline
        $f'+g'$ & $f+g$ & $f$, $g$ dérivables \\
        \hline
        $\lambda f'$ & $\lambda f$ & $f$ dérivable \\
        \hline
        $f'\,f^n$, $n\in\mathbb{Z}$, $n\neq-1$ & $\dfrac{f^{n+1}}{n+1}$ & $f$ dérivable \\
        \hline
        $\dfrac{f'}{f^2}$ & $-\dfrac{1}{f}$ & $f$ dérivable, $f\neq 0$ \\
        \hline
        $\dfrac{f'}{\sqrt{f}}$ & $2\sqrt{f}$ & $f$ dérivable, $f>0$ \\
        \hline
        $\dfrac{f'}{f}$ & $\ln|f|$ & $f$ dérivable, $f\neq 0$ \\
        \hline
        $f'\,e^f$ & $e^f$ & $f$ dérivable \\
        \hline
        $f'\times(g'\circ f)$ & $g\circ f$ & \\
        \hline
    \end{tabular}

    \section{Équation différentielle $y'=ay$}
    \begin{Propriété}
        Les solutions de $y'=ay$ ($a\neq 0$) sont les fonctions
        $x\mapsto ke^{ax}$, $k\in\mathbf{R}$.
        Pour tout $(x_0,y_0)$, il existe une unique solution vérifiant $y(x_0)=y_0$.
    \end{Propriété}
    \begin{proof}
        $f_k(x)=ke^{ax}$ vérifie $f_k'(x)=kae^{ax}=af_k(x)$, donc $f_k$ est solution.\\
        Réciproquement, si $y$ est solution, posons $g(x)=y(x)e^{-ax}$.
        Alors $g'(x)=y'(x)e^{-ax}-ay(x)e^{-ax}=ay(x)e^{-ax}-ay(x)e^{-ax}=0$,
        donc $g$ est constante : $g(x)=k$, soit $y(x)=ke^{ax}$.
    \end{proof}

    \section{Équation différentielle $y'=ay+b$}
    \begin{Propriété}
        Les solutions de $y'=ay+b$ ($a,b\neq 0$) sont les fonctions
        $x\mapsto ke^{ax}-\dfrac{b}{a}$, $k\in\mathbf{R}$.
    \end{Propriété}
    \begin{proof}
        La solution constante particulière vérifie $0=ak+b$, soit $k=-\dfrac{b}{a}$.
        Si $g$ est une solution quelconque et $f(x)=-\dfrac{b}{a}$, alors $h=g-f$
        vérifie $h'=a h$, donc $h(x)=ke^{ax}$, d'où $g(x)=ke^{ax}-\dfrac{b}{a}$.
    \end{proof}

    \section{Équation différentielle $y'=ay+\varphi$}
    \begin{Propriété}
        Si $g$ est une solution particulière de $y'=ay+\varphi$, les solutions générales
        sont $x\mapsto ke^{ax}+g(x)$, $k\in\mathbf{R}$.
    \end{Propriété}

    \chapter{Combinatoire et dénombrement}
    \nchap

    \begin{capacitesbox}
    \begin{itemize}
        \item Utiliser les règles du dénombrement : principe additif et multiplicatif.
        \item Calculer le nombre de permutations de $n$ éléments : $n!$
        \item Calculer le nombre d'arrangements de $k$ éléments parmi $n$ : $A_n^k=\frac{n!}{(n-k)!}$.
        \item Calculer les coefficients binomiaux $\binom{n}{k}$ et utiliser la formule de Pascal.
        \item Développer $(a+b)^n$ à l'aide du binôme de Newton.
    \end{itemize}
    \end{capacitesbox}

    \begin{demoprogbox}
    \begin{itemize}
        \item Démontrer que le nombre de parties d'un ensemble à $n$ éléments est $2^n$
              (bijection avec les mots binaires de longueur $n$).
        \item Démontrer la formule de Pascal :
              $\dbinom{n}{k}+\dbinom{n}{k+1}=\dbinom{n+1}{k+1}$.
    \end{itemize}
    \end{demoprogbox}

    \section{Arrangements et permutations}
    \begin{Définition}
        On appelle \textbf{factorielle} de $n$ le nombre noté :
        \[n!=n\times(n-1)\times\cdots\times 2\times 1\]
        Par convention $0!=1$.
    \end{Définition}
    \begin{Définition}
        Soient $A$ un ensemble fini non vide à $n$ éléments et $k$ un entier naturel
        inférieur ou égal à $n$. Un \textbf{arrangement} de $k$ éléments de $A$ (ou
        $k$-arrangement) est un $k$-uplet d'éléments distincts de $A$.
    \end{Définition}
    \begin{Propriété}
        Soient $n$ un entier naturel non nul et $k$ un entier compris entre $1$ et $n$.
        Le nombre de $k$-arrangements d'un ensemble $E$ à $n$ éléments est :
        \[n\times(n-1)\times\cdots\times(n-k+1)=\frac{n!}{(n-k)!}.\]
    \end{Propriété}
    \begin{Définition}
        Une \textbf{permutation} d'un ensemble $E$ à $n$ éléments est un $n$-uplet
        d'éléments distincts de $E$.
    \end{Définition}
    \begin{Propriété}
        Le nombre de permutations d'un ensemble $E$ à $n$ éléments est $n!$.
    \end{Propriété}
    \section{Combinaisons}
    \begin{Propriété}
        Le nombre de parties d'un ensemble fini à $n$ éléments est $2^n$.
    \end{Propriété}
    \begin{proof}
        Chaque élément peut être inclus (1) ou exclu (0) d'une partie.
        Il y a donc $2^n$ suites binaires de longueur $n$, correspondant aux $2^n$ parties.
    \end{proof}
    \subsection{Nombre de combinaisons}
    \begin{Définition}
        Une \textbf{combinaison} de $k$ éléments de $E$ est une partie de $E$ à $k$
        éléments. Le nombre de combinaisons de $k$ éléments parmi $n$ est noté
        $\dbinom{n}{k}$.
    \end{Définition}
    \begin{Propriété}
        Pour $k\leq n$ :
        \begin{enumerate}
            \item $\dbinom{n}{k}=\dfrac{n!}{k!\,(n-k)!}$ et $\dbinom{n}{k}=\dbinom{n}{n-k}$.
            \item \textbf{Relation de Pascal :} si $1\leq k\leq n-1$,
                  $\dbinom{n}{k}=\dbinom{n-1}{k-1}+\dbinom{n-1}{k}$.
            \item $\dbinom{n}{0}=1$, $\dbinom{n}{1}=n$,
                  $\dbinom{n}{2}=\dfrac{n(n-1)}{2}$.
        \end{enumerate}
    \end{Propriété}

    \chapter{Fonctions cosinus et sinus}
    \nchap

    \begin{capacitesbox}
    \begin{itemize}
        \item Connaître les valeurs exactes de $\cos$ et $\sin$ pour les angles remarquables.
        \item Utiliser les formules de somme $\cos(a\pm b)$, $\sin(a\pm b)$ et de double-angle.
        \item Résoudre les équations $\cos x=k$, $\sin x=k$, $\cos x=\cos a$, $\sin x=\sin a$.
        \item Étudier les fonctions faisant intervenir $\cos$ ou $\sin$ (dérivée, tableau de variations).
    \end{itemize}
    \end{capacitesbox}

    \begin{demoprogbox}
    \begin{itemize}
        \item Démontrer les formules d'addition $\cos(a+b)=\cos a\cos b-\sin a\sin b$
              et $\sin(a+b)=\sin a\cos b+\cos a\sin b$ (à partir de la rotation ou du produit scalaire).
        \item Démontrer que $(\cos x)'=-\sin x$ et $(\sin x)'=\cos x$.
        \item Retrouver $\displaystyle\lim_{x\to 0}\dfrac{\sin x}{x}=1$ à partir des dérivées ou d'un encadrement.
    \end{itemize}
    \end{demoprogbox}

    \section{Fonctions sinus et cosinus}
    \subsection{Généralités}
    \begin{Définition}
        Dans un repère orthogonal, soient $x\in\mathbf{R}$ et $M$ le point du cercle
        trigonométrique associé à $x$.
        \begin{itemize}
            \item $\sin(x)$ est l'ordonnée de $M$.
            \item $\cos(x)$ est l'abscisse de $M$.
        \end{itemize}
    \end{Définition}

    \begin{Propriété}
        \begin{itemize}
            \item La fonction cosinus est \textbf{paire} et $2\pi$-périodique.
            \item La fonction sinus est \textbf{impaire} et $2\pi$-périodique.
        \end{itemize}
    \end{Propriété}

    \begin{Propriété}
        Pour tout $a\in\mathbf{R}$ et $x\in\mathbf{R}$ :
        \begin{itemize}
            \item $\cos x=\cos a \Longleftrightarrow x=a+2k\pi$ ou $x=-a+2k\pi$,
                  $k\in\mathbf{Z}$.
            \item $\sin x=\sin a \Longleftrightarrow x=a+2k\pi$ ou $x=\pi-a+2k\pi$,
                  $k\in\mathbf{Z}$.
        \end{itemize}
    \end{Propriété}

    \subsection{Dérivées}
    \begin{Propriété}
        $\cos'(x)=-\sin(x)$ et $\sin'(x)=\cos(x)$.
    \end{Propriété}

    \begin{Propriété}
        Si $u$ est dérivable sur $I$ :
        $(\cos\circ\, u)'(x)=-u'(x)\sin(u(x))$ et
        $(\sin\circ\, u)'(x)=u'(x)\cos(u(x))$.
    \end{Propriété}

    \subsection{Primitives}
    \begin{Propriété}
        \begin{itemize}
            \item $u'(x)\sin(u(x))$ a pour primitive $-\cos(u(x))$.
            \item $u'(x)\cos(u(x))$ a pour primitive $\sin(u(x))$.
        \end{itemize}
    \end{Propriété}

    \section*{Tableaux de variation et représentations graphiques}

    \subsection*{Fonction cosinus sur $[0;\pi]$}
    \begin{center}
    \begin{tikzpicture}
    \tkzTabInit[lgt=1.5,espcl=2.5]{$x$/1, $\cos(x)$/2}{$0$, $\pi$}
    \tkzTabVar{+/$1$, -/$-1$}
    \end{tikzpicture}
    \end{center}

    \begin{center}
    \begin{pspicture}(-4,-2)(4,2)
        \psaxes{->}(0,0)(-4,-1.5)(4,1.5)
        \psplot[algebraic=true]{-4}{4}{cos(x)}
    \end{pspicture}
    \end{center}

    \subsection*{Fonction sinus sur $[0;\pi]$}
    La fonction sinus est croissante sur $[0;\frac{\pi}{2}]$ et décroissante sur
    $[\frac{\pi}{2};\pi]$.

    \begin{center}
    \begin{tikzpicture}
    \tkzTabInit[lgt=1.5,espcl=2.5]{$x$/1, $\sin(x)$/2}{$0$, $\dfrac{\pi}{2}$, $\pi$}
    \tkzTabVar{-/$0$, +/$1$, -/$0$}
    \end{tikzpicture}
    \end{center}

    \begin{center}
    \begin{pspicture}(-4,-2)(4,2)
        \psaxes{->}(0,0)(-4,-1.5)(4,1.5)
        \psplot[algebraic=true]{-4}{4}{sin(x)}
    \end{pspicture}
    \end{center}

    \chapter{Calcul intégral}
    \nchap

    \begin{capacitesbox}
    \begin{itemize}
        \item Calculer une intégrale à l'aide d'une primitive : $\displaystyle\int_a^b f(x)\,\mathrm{d}x = [F(x)]_a^b$.
        \item Interpréter l'intégrale comme une aire algébrique.
        \item Calculer une valeur moyenne d'une fonction sur $[a;b]$.
        \item Utiliser la linéarité, la relation de Chasles et la positivité de l'intégrale.
        \item Effectuer une intégration par parties : $\displaystyle\int_a^b uv'\,\mathrm{d}x=[uv]_a^b-\int_a^b u'v\,\mathrm{d}x$.
    \end{itemize}
    \end{capacitesbox}

    \begin{demoprogbox}
    \begin{itemize}
        \item Démontrer que si $f\geq 0$ sur $[a;b]$, alors $\displaystyle\int_a^b f(x)\,\mathrm{d}x\geq 0$.
        \item Démontrer l'inégalité de la moyenne :
              $\displaystyle\left|\int_a^b f(x)\,\mathrm{d}x\right|\leq(b-a)\max_{[a;b]}|f|$.
        \item Démontrer la formule d'intégration par parties.
    \end{itemize}
    \end{demoprogbox}

    \section{Intégrale et primitive}
    \subsection{Intégrale d'une fonction positive}
    \begin{Définition}
        Soit $f$ continue et positive sur $[a;b]$. L'\textbf{intégrale} de $a$ à $b$ de $f$
        est l'aire (en unités d'aire) du domaine délimité par la courbe, l'axe des abscisses
        et les droites $x=a$ et $x=b$ :
        \[
            \int_a^b f(x)\,\mathrm{d}x.
        \]
    \end{Définition}

    \begin{Théorème}[Existence d'une primitive]
        Soit $f$ continue et positive sur $[a;b]$. La fonction $F$ définie par
        $F(x)=\int_a^x f(t)\,\mathrm{d}t$ est dérivable sur $[a;b]$ et $F'=f$.
        C'est la primitive de $f$ qui s'annule en $a$.
    \end{Théorème}

    \begin{Propriété}
        Si $f$ est continue et négative sur $[a;b]$, l'aire du domaine est
        $\displaystyle\int_a^b -f(x)\,\mathrm{d}x$.
    \end{Propriété}

    \begin{Propriété}
        Si $F$ est une primitive quelconque de $f$ sur $[a;b]$ :
        \[
            \int_a^b f(t)\,\mathrm{d}t = F(b)-F(a) = \bigl[F(x)\bigr]_a^b.
        \]
    \end{Propriété}

    \subsection{Intégrale d'une fonction de signe quelconque}
    \begin{Propriété}
        Soient $f$ et $g$ continues sur $I$, et $a,b,c\in I$, $\lambda\in\mathbf{R}$.
        \begin{itemize}
            \item $\int_a^a f(x)\,\mathrm{d}x=0$
            \item $\int_b^a f(x)\,\mathrm{d}x=-\int_a^b f(x)\,\mathrm{d}x$
            \item \textbf{Relation de Chasles :}
                  $\int_a^c f = \int_a^b f + \int_b^c f$
            \item \textbf{Linéarité :}
                  $\int_a^b(f+g) = \int_a^b f + \int_a^b g$ et
                  $\int_a^b\lambda f = \lambda\int_a^b f$
        \end{itemize}
    \end{Propriété}

    \subsection{Signe de l'intégrale}
    \begin{Propriété}
        Pour $a\leq b$ et $f$ continue sur $[a;b]$ :
        \begin{enumerate}
            \item Si $f(x)\geq 0$ sur $[a;b]$, alors $\int_a^b f(x)\,\mathrm{d}x\geq 0$.
            \item Si $f(x)\leq 0$ sur $[a;b]$, alors $\int_a^b f(x)\,\mathrm{d}x\leq 0$.
            \item Si $f(x)\geq g(x)$ sur $[a;b]$, alors
                  $\int_a^b f\,\mathrm{d}x\geq\int_a^b g\,\mathrm{d}x$.
        \end{enumerate}
    \end{Propriété}

    \subsection{Intégration par parties}
    \begin{Théorème}
        Si $u$ et $v$ sont dérivables sur $[a;b]$ avec $u'$ et $v'$ continues :
        \[
            \int_a^b u(x)v'(x)\,\mathrm{d}x
            = \bigl[u(x)v(x)\bigr]_a^b - \int_a^b u'(x)v(x)\,\mathrm{d}x.
        \]
    \end{Théorème}

    \section{Applications du calcul intégral}
    \subsection{Aire entre deux courbes}
    \begin{Propriété}
        Si $f(x)\leq g(x)$ sur $[a;b]$, l'aire du domaine compris entre les courbes de $f$
        et de $g$ est :
        \[
            \int_a^b\bigl(g(x)-f(x)\bigr)\,\mathrm{d}x.
        \]
    \end{Propriété}

    \subsection{Valeur moyenne d'une fonction}
    \begin{Propriété}
        La \textbf{valeur moyenne} de $f$ continue sur $[a;b]$ est :
        \[
            \mu = \frac{1}{b-a}\int_a^b f(x)\,\mathrm{d}x.
        \]
        Graphiquement, $\mu(b-a)=\int_a^b f(x)\,\mathrm{d}x$ : l'aire du rectangle de
        hauteur $\mu$ égale l'aire sous la courbe.
    \end{Propriété}

    \chapter{Loi des grands nombres}
    \nchap

    \begin{capacitesbox}
    \begin{itemize}
        \item Calculer l'espérance, la variance et l'écart-type d'une variable aléatoire.
        \item Comprendre et utiliser l'inégalité de Bienaymé-Tchebychev.
        \item Interpréter et appliquer la loi des grands nombres (convergence de la fréquence).
        \item Utiliser la concentration autour de la moyenne pour prendre des décisions.
    \end{itemize}
    \end{capacitesbox}

    \begin{demoprogbox}
    \begin{itemize}
        \item Démontrer $E(X)=np$ et $V(X)=np(1-p)$ pour $X\sim\mathcal{B}(n;p)$.
        \item Démontrer l'inégalité de Bienaymé-Tchebychev :
              $P(|X-E(X)|\geq\delta)\leq\dfrac{V(X)}{\delta^2}$.
        \item Démontrer la loi faible des grands nombres (admise ; justifier le cas $p=\frac{1}{2}$).
    \end{itemize}
    \end{demoprogbox}

    \begin{Définition}
        Soit $X$ une variable aléatoire. La variable $Y=aX+b$ prend les valeurs
        $y_i=ax_i+b$.
    \end{Définition}

    \begin{Propriété}
        Pour $Y=aX+b$ :
        \begin{itemize}
            \item $E(Y)=aE(X)+b$
            \item $V(Y)=a^2V(X)$
            \item $\sigma(aX)=|a|\sigma(X)$
        \end{itemize}
    \end{Propriété}

    \section{Inégalités}
    \subsection{Inégalité de Markov}
    \begin{Propriété}
        Soit $X$ une variable aléatoire à valeurs positives et $a>0$ :
        \[
            P(X\geq a)\leq\frac{E(X)}{a}.
        \]
    \end{Propriété}

    \subsection{Inégalité de Bienaymé-Tchebychev}
    \begin{Propriété}
        Soit $X$ une variable aléatoire et $a>0$ :
        \[
            P\!\left(|X-E(X)|\geq a\right)\leq\frac{V(X)}{a^2}.
        \]
    \end{Propriété}

    \section{Inégalité de concentration}
    \begin{Définition}
        Un \textbf{échantillon de taille $n$} de la loi de $X$ est une liste
        $(X_1,\ldots,X_n)$ de variables aléatoires indépendantes de même loi que $X$.
    \end{Définition}

    \begin{Théorème}
        Soit $M_n=\dfrac{1}{n}\sum_{i=1}^n X_i$. Pour tout $a>0$ :
        \[
            P\!\left(|M_n-E(X)|\geq a\right)\leq\frac{V(X)}{na^2}.
        \]
    \end{Théorème}

    \begin{Propriété}[Loi des grands nombres]
        Pour tout $a>0$ :
        \[
            \lim_{n\to+\infty}P\!\left(|M_n-E(X)|\geq a\right)=0.
        \]
        La moyenne empirique converge en probabilité vers l'espérance.
    \end{Propriété}


    \appendix
    \chapter{Démonstrations}
    \nchap

    \section{Algèbre et combinatoire}

    \subsection{Identité $\sum_{k=0}^{n} \binom{n}{k} = 2^n$}

    \begin{enoncebox}
    Soit $n \in \mathbb{N}$. On veut démontrer l'identité :
    \[\sum_{k=0}^{n} \binom{n}{k} = 2^n.\]
    \end{enoncebox}

    \begin{methodebox}{Méthode 1 — Binôme de Newton}
    La formule du binôme de Newton stipule que pour tous réels $a$, $b$ et tout $n\in\mathbf{N}$ :
    \[(a + b)^n = \sum_{k=0}^{n} \binom{n}{k} a^k b^{n-k}.\]
    Pour $a = b = 1$ :
    \begin{align*}
    2^n = (1+1)^n &= \sum_{k=0}^{n} \binom{n}{k} 1^k \cdot 1^{n-k} = \sum_{k=0}^{n} \binom{n}{k}.
    \end{align*}
    \end{methodebox}

    \begin{methodebox}{Méthode 2 — Dénombrement}
    Soit $E$ un ensemble de cardinal $n$. On compte les parties de $E$ de deux façons.

    \textbf{Première approche :} pour chaque élément, on choisit de l'inclure (1) ou non (0), soit $2^n$ parties.

    \textbf{Deuxième approche :} on classe les parties par taille $k$ : il y en a $\binom{n}{k}$ à $k$ éléments. En sommant :
    $\displaystyle\sum_{k=0}^{n}\binom{n}{k}$ parties au total.

    Ces deux expressions sont égales, d'où $\displaystyle\sum_{k=0}^{n}\binom{n}{k}=2^n$.
    \end{methodebox}

    \subsection{Relation de Pascal}

    \begin{enoncebox}
    Pour tout entier $n$ et tout $k$ tel que $1\leq k\leq n$ :
    \[\binom{n}{k}+\binom{n}{k-1}=\binom{n+1}{k}.\]
    \end{enoncebox}

    \begin{methodebox}{Méthode 1 — Calcul factoriel}
    \begin{align*}
    \binom{n}{k}+\binom{n}{k-1}
    &=\frac{n!}{k!(n-k)!}+\frac{n!}{(k-1)!(n-k+1)!}\\
    &=\frac{n!(n-k+1)}{k!(n-k+1)!}+\frac{n!\,k}{k!(n-k+1)!}\\
    &=\frac{n!(n-k+1+k)}{k!(n-k+1)!}
    =\frac{(n+1)!}{k!(n+1-k)!}=\binom{n+1}{k}.
    \end{align*}
    \end{methodebox}

    \begin{methodebox}{Méthode 2 — Interprétation combinatoire}
    Soit $E$ un ensemble de $n+1$ éléments et $x\in E$ un élément fixé.
    Tout sous-ensemble de $k$ éléments de $E$ soit contient $x$ (il reste à choisir $k-1$ parmi les $n$ autres : $\binom{n}{k-1}$ façons), soit ne contient pas $x$ ($\binom{n}{k}$ façons).
    Ces cas sont disjoints et exhaustifs, donc $\binom{n+1}{k}=\binom{n}{k-1}+\binom{n}{k}$.
    \end{methodebox}

    \subsection{Équation cartésienne d'un plan}

    \begin{enoncebox}
    Soient $A(x_A,y_A,z_A)$ et $\vc{n}\vcb{a}{b}{c}$ un vecteur non nul.
    L'équation du plan $\mathcal{P}$ passant par $A$ de vecteur normal $\vc{n}$ est de la forme $ax+by+cz+d=0$.
    \end{enoncebox}

    \begin{methodebox}{Démonstration — Produit scalaire}
    $M(x,y,z)\in\mathcal{P}$ si et seulement si $\vc{AM}\perp\vc{n}$, soit $\vc{AM}\cdot\vc{n}=0$.
    Or $\vc{AM}=\vcb{x-x_A}{y-y_A}{z-z_A}$, donc :
    \begin{align*}
    a(x-x_A)+b(y-y_A)+c(z-z_A)&=0\\
    ax+by+cz+d&=0
    \end{align*}
    avec $d=-(ax_A+by_A+cz_A)$.
    \end{methodebox}

    \section{Probabilités}

    \subsection{Loi binomiale : probabilité de $k$ succès}

    \begin{enoncebox}
    Pour $X\sim\mathcal{B}(n;p)$ et $k\in\{0,\ldots,n\}$ :
    $P(X=k)=\dbinom{n}{k}p^k(1-p)^{n-k}$.
    \end{enoncebox}

    \begin{methodebox}{Démonstration — Schéma de Bernoulli}
    Un chemin comportant $k$ succès et $n-k$ échecs a probabilité $p^k(1-p)^{n-k}$ (indépendance).
    Le nombre de tels chemins est $\binom{n}{k}$ (choix des positions des succès).
    Ces chemins étant incompatibles : $P(X=k)=\binom{n}{k}p^k(1-p)^{n-k}$.
    \end{methodebox}

    \subsection{Espérance de la loi binomiale}

    \begin{methodebox}{Méthode 1 — Linéarité de l'espérance}
    On écrit $X=X_1+\cdots+X_n$ où chaque $X_i\sim\mathcal{B}(p)$.
    Pour chaque variable de Bernoulli : $E(X_i)=p$.
    Par linéarité : $E(X)=\displaystyle\sum_{i=1}^n E(X_i)=np$.
    \end{methodebox}

    \begin{methodebox}{Méthode 2 — Définition de l'espérance}
    $E(X)=\displaystyle\sum_{k=0}^{n}k\binom{n}{k}p^k(1-p)^{n-k}$.
    Pour $k\geq 1$, $k\binom{n}{k}=n\binom{n-1}{k-1}$, d'où :
    \[E(X)=np\sum_{j=0}^{n-1}\binom{n-1}{j}p^j(1-p)^{n-1-j}=np\cdot(p+(1-p))^{n-1}=np.\]
    \end{methodebox}

    \subsection{Variance de la loi binomiale}

    \begin{methodebox}{Méthode 1 — Variables indépendantes}
    $V(X_i)=E(X_i^2)-[E(X_i)]^2=p-p^2=p(1-p)$.
    Par indépendance : $V(X)=\displaystyle\sum_{i=1}^nV(X_i)=np(1-p)$.
    \end{methodebox}

    \begin{methodebox}{Méthode 2 — Définition de la variance}
    On utilise $V(X)=E(X^2)-[E(X)]^2$.
    En écrivant $k^2=k(k-1)+k$ et $k(k-1)\binom{n}{k}=n(n-1)\binom{n-2}{k-2}$ :
    \[E(X^2)=n(n-1)p^2+np,\]
    donc $V(X)=n(n-1)p^2+np-(np)^2=np(1-p)$.
    \end{methodebox}

    \section{Analyse}

    \subsection{Limite d'une suite croissante non majorée}

    \begin{enoncebox}
    Si $(u_n)$ est croissante et non majorée, alors $\lim_{n\to+\infty}u_n=+\infty$.
    \end{enoncebox}

    \begin{methodebox}{Démonstration}
    Soit $A\in\mathbf{R}$. Comme $(u_n)$ n'est pas majorée, $A$ n'est pas un majorant : il existe $N$ tel que $u_N>A$.
    Comme $(u_n)$ est croissante, pour tout $n\geq N$ : $u_n\geq u_N>A$.
    Donc pour tout $A$, à partir d'un certain rang tous les termes dépassent $A$ : $\lim u_n=+\infty$.
    \end{methodebox}

    \subsection{Inégalité de Bernoulli}

    \begin{enoncebox}
    Pour tout $a>0$ et tout $n\in\mathbf{N}$ : $(1+a)^n\geq 1+na$.
    \end{enoncebox}

    \begin{methodebox}{Démonstration par récurrence}
    \textbf{Initialisation :} pour $n=0$, $(1+a)^0=1=1+0\cdot a$. Vrai.

    \textbf{Hérédité :} supposons $(1+a)^k\geq 1+ka$. Comme $1+a>0$ :
    \begin{align*}
    (1+a)^{k+1}&=(1+a)^k(1+a)\geq(1+ka)(1+a)\\
    &=1+(k+1)a+ka^2\geq 1+(k+1)a.
    \end{align*}
    \textbf{Conclusion :} la propriété est vraie pour tout $n\in\mathbf{N}$.
    \end{methodebox}

    \subsection{Limite de la suite $(q^n)$}

    \begin{methodebox}{Cas $q>1$}
    Posons $q=1+a$ avec $a>0$. Par l'inégalité de Bernoulli : $q^n=(1+a)^n\geq 1+na\to+\infty$.
    Donc $\lim q^n=+\infty$.
    \end{methodebox}

    \begin{methodebox}{Cas $0<q<1$}
    Posons $Q=\frac{1}{q}>1$. D'après le cas précédent, $Q^n\to+\infty$, donc $q^n=\frac{1}{Q^n}\to 0$.
    \end{methodebox}

    \subsection{Convexité : la courbe est au-dessus de ses tangentes}

    \begin{enoncebox}
    Si $f''(x)\geq 0$ sur $I$, alors $\mathcal{C}_f$ est au-dessus de chacune de ses tangentes.
    \end{enoncebox}

    \begin{methodebox}{Démonstration}
    Soit $a\in I$. La tangente en $a$ a pour équation $y=f'(a)(x-a)+f(a)$.
    Posons $g(x)=f(x)-f'(a)(x-a)-f(a)$. Alors $g(a)=0$ et $g'(x)=f'(x)-f'(a)$.
    Comme $f''\geq 0$, $f'$ est croissante : $g'(x)\leq 0$ pour $x\leq a$ et $g'(x)\geq 0$ pour $x\geq a$.
    Donc $g$ admet un minimum en $a$ valant $0$ : pour tout $x\in I$, $g(x)\geq 0$, soit $f(x)\geq f'(a)(x-a)+f(a)$.
    \end{methodebox}

    \subsection{Croissances comparées : $\lim_{x\to+\infty}\frac{e^x}{x}=+\infty$}

    \begin{enoncebox}
    Pour tout $n\in\mathbf{N}$ : $\lim_{x\to+\infty}\frac{e^x}{x^n}=+\infty$.
    En particulier, $\lim_{x\to+\infty}\frac{e^x}{x}=+\infty$.
    \end{enoncebox}

    \begin{methodebox}{Démonstration --- cas n = 1}
    De la convexité de l'exponentielle, on tire $e^u\geq u+1>u$ pour tout $u\in\mathbf{R}$.
    Posons $u=\frac{x}{2}$ : $e^{x/2}>\frac{x}{2}$. En élevant au carré : $e^x>\frac{x^2}{4}$.
    Pour $x>0$ : $\frac{e^x}{x}>\frac{x}{4}\to+\infty$.
    Par comparaison : $\lim_{x\to+\infty}\frac{e^x}{x}=+\infty$.
    \end{methodebox}

    \subsection{Croissance comparée : $\lim_{x\to-\infty}x\,e^x=0$}

    \begin{enoncebox}
    $\lim_{x\to-\infty}x\,e^x=0$.
    \end{enoncebox}

    \begin{methodebox}{Démonstration}
    Posons $X=-x$ ; quand $x\to-\infty$, $X\to+\infty$.
    \[x\,e^x=(-X)e^{-X}=-\frac{X}{e^X}=-\frac{1}{\frac{e^X}{X}}.\]
    Or $\lim_{X\to+\infty}\frac{e^X}{X}=+\infty$, donc $\lim_{x\to-\infty}x\,e^x=0$.
    \end{methodebox}

    \subsection{Limite de $x\ln x$ en $0^+$}

    \begin{enoncebox}
    $\lim_{x\to 0^+}x\ln x=0$.
    \end{enoncebox}

    \begin{methodebox}{Démonstration}
    Posons $x=e^t$ : quand $x\to 0^+$, $t\to-\infty$ et $x\ln x=e^t\cdot t=t\,e^t$.
    Or $\lim_{t\to-\infty}t\,e^t=0$ (démonstration précédente), donc $\lim_{x\to 0^+}x\ln x=0$.
    \end{methodebox}

    \subsection{Deux primitives d'une même fonction diffèrent d'une constante}

    \begin{enoncebox}
    Soient $F$ et $G$ deux primitives de $f$ sur $I$. Il existe $C\in\mathbf{R}$ tel que $G(x)=F(x)+C$.
    \end{enoncebox}

    \begin{methodebox}{Démonstration}
    Posons $H=G-F$. Alors $H'(x)=G'(x)-F'(x)=f(x)-f(x)=0$ pour tout $x\in I$.
    Une fonction de dérivée nulle sur un intervalle est constante : $H(x)=C$.
    Donc $G(x)=F(x)+C$.
    \end{methodebox}

\end{document}
 